\documentclass[10pt]{amsart}
\usepackage[a4paper,margin=23mm]{geometry}
\usepackage{amsmath,amssymb,amsthm,mathtools,hyperref,graphicx,xspace,letltxmacro,xpatch}
\usepackage[english]{babel}
\usepackage{xfrac,xcolor}
\newtheorem{cdrthm}{Theorem}[section]
\newtheorem{theo}[cdrthm]{Theorem}
\newtheorem{prop}[cdrthm]{Proposition}
\newtheorem{lemm}[cdrthm]{Lemma}
\newtheorem{coro}[cdrthm]{Corollary}
\theoremstyle{definition}
\newtheorem{rema}[cdrthm]{Remark}
\newenvironment{proofc}{\begin{proof}[Proof (continued)]}{\end{proof}}


\newtheorem{defi}[cdrthm]{Definition}
\newtheorem{lemma}[cdrthm]{Lemma}
\newtheorem*{conj*}{Conjecture}
\newtheorem{sourceenonce}[cdrthm]{\SourceEnonceName}
\NewDocumentEnvironment{enonce}{O{plain} m o}{\def\SourceEnonceName{#2}\IfNoValueTF{#3}{\begin{sourceenonce}}{\begin{sourceenonce}[{#3}]}}{\end{sourceenonce}}

\usepackage{bm}




\newcommand{\TF}{\mathrm{TF}}
\DeclareMathOperator{\sgn}{sgn}
\DeclareMathOperator{\tr}{tr}
\DeclareMathOperator{\const}{const}
\newcommand{\loc}{\mathrm{loc}}
\newcommand{\spanl}{\mathrm{span}}


\newcommand*\rd{\mathrm{d}}
\newcommand*\ri{\mathrm{i}}
\newcommand{\ess}{\mathrm{ess}}
\newcommand{\pp}{\mathrm{pp}}


\newcommand*\dr{\mathrm{d}r}
\newcommand*\dx{\mathrm{d}x}
\newcommand*\dy{\mathrm{d}y}



\newcommand{\bbS}{\mathbb{S}}
\newcommand*\cz{\mathbb{C}}
\newcommand*\nz{\mathbb{N}}
\newcommand*\gz{\mathbb{Z}}
\newcommand*\rz{\mathbb{R}}


\newcommand*\zp{\dot{\gz}}

%\newcommand*\bbS{\bbS}
\newcommand{\one}{\mathbf{1}}

\newcommand*\cD{\mathcal{D}}
\newcommand*\cE{\mathcal{E}}
\newcommand*\cK{\mathcal{K}}
\newcommand*\cS{\mathcal{S}}
\newcommand*\cO{\mathcal{O}}


\newcommand*\gD{\mathfrak{D}}
\newcommand*\gH{\mathfrak{H}}
\newcommand*\gh{\mathfrak{h}}
\newcommand*\gQ{\mathfrak{Q}}
\newcommand*\gS{\mathfrak{S}}

\newcommand*\ghk{{\gh_{\kappa}}}
\newcommand*\ghpk{{\gh^+_\kappa}}
\newcommand*\ppp{\Pi^+_\kappa}


%\newcommand*\const{\mathrm{const}\,}
%\gS

\newcommand*\domega{\rd\omega}


\newcommand*\balpha{\bm{\alpha}}
\newcommand*\bsigma{\bm{\sigma}}
\newcommand*\bp{\bm{p}}

%%%%%%%%%%%%%%%%%%%%%%%%%%%%%%%%%%%%%%%%%%%%%%%%%%%%%%%%%%%%%%%%%%%%%%%


\graphicspath{{./figures/}}

\newcommand*{\mk}{\mkern -1mu}
\newcommand*{\Mk}{\mkern -2mu}
\newcommand*{\mK}{\mkern 1mu}
\newcommand*{\MK}{\mkern 2mu}

\hypersetup{urlcolor=purple, linkcolor=blue, citecolor=red}

\newcommand*{\romanenumi}{\renewcommand*{\theenumi}{\roman{enumi}}}
\newcommand*{\Romanenumi}{\renewcommand*{\theenumi}{\Roman{enumi}}}
\newcommand*{\alphenumi}{\renewcommand*{\theenumi}{\alph{enumi}}}
\newcommand*{\Alphenumi}{\renewcommand*{\theenumi}{\Alph{enumi}}}
\let\oldtilde\tilde
\renewcommand*{\tilde}[1]{\mathchoice{\widetilde{#1}}{\widetilde{#1}}{\oldtilde{#1}}{\oldtilde{#1}}}
%%%%%%%%%%%%%%%%%%%%%%%%%%%%%%%%%%%%%%%%%%%%%%%%%%%%%%%%%%%%%%%%%

\let\dpl\displaystyle
\providecommand{\skpt}{\par\smallskip\noindent}
\providecommand{\pointir}{.\ }
\providecommand{\eg}{e.g.}\providecommand{\ie}{i.e.}\providecommand{\etc}{etc.}\providecommand{\resp}{resp.}\providecommand{\cf}{cf.}
\providecommand{\mainmatter}{}\providecommand{\backmatter}{}
\let\Subsection\subsection\let\Subsubsection\subsubsection
\setlength{\emergencystretch}{2em}

\AtBeginDocument{\gdef\shorttitle{Strong Scott density limit in the Furry picture}}
\title{Strong Scott density convergence in the relativistic Furry picture}
\author{Konstantin Merz, Heinz Siedentop}
\date{doi:10.5802/ahl.131; CC BY 4.0}
\begin{document}
\maketitle
\noindent\textit{Editorial note.} Complete original mathematical body and bibliography retained. Authors retain responsibility for mathematics.\par\medskip
\noindent\textit{Editorial note.} Original neutral Furry model, fixed gamma=Z/c in(0,1) and the exact radial test-function spaces are preserved. No pointwise-density convergence, delta tests or critical gamma endpoint is added.\par\medskip
\noindent\textit{Allowed standard prerequisites (editorial).}\begin{itemize}
\item Author-cited relativistic Hardy estimates, Scott energy asymptotics and spectral trace/Feynman-Hellmann facts. Original cited prior inputs remain cited; all source-local constructions and complete proofs are retained.
\end{itemize}\medskip
\section{Introduction}
The quest for ground state properties of Coulomb systems like atoms, molecules, and solids is one of the central topics in physics and chemistry. However, it became clear right after the discovery of quantum mechanics (Heisenberg~\cite{Heisenberg1925}) that -- not much different from classical mechanics -- one-particle problems like the hydrogen atom can be solved analytically (Pauli~\cite{Pauli1926}) but problems with several electrons need suitable approximations. Within two years after the advent of quantum mechanics Thomas~\cite{Thomas1927} and Fermi~\cite{Fermi1927,Fermi1928}) developed an approximation-- now called Thomas--Fermi theory -- for predicting the ground state energies and densities of large atoms. About fifty years later Lieb and Simon~\cite{LiebSimon1977} showed in their seminal work that, indeed, the asymptotic behavior of atomic energies for large atomic numbers $Z$ is given by the Thomas--Fermi energy, and that the suitably renormalized ground state density of large atoms on the scale $Z^{-1/3}$ converges to the hydrogenic Thomas--Fermi density.

The Thomas--Fermi theory is the simplest example of what is called density functional theory. However, already the next order correction is not easily connected with the first correction -- given by the Thomas--Fermi--Weizs\"acker functional~\cite{Weizsacker1935} -- of Thomas--Fermi theory. It requires a renormalization of the constant in front of the inhomogeneity correction (Yonei and Tomishima~\cite{YoneiTomishima1965}, see also Lieb~\cite{Lieb1982A,LiebLiberman1982}). In fact the next order energy correction was predicted by Scott~\cite{Scott1952} as stemming entirely from the electrons on the scale $Z^{-1}$ where the interaction between the electrons is completely dominated by the electron-nucleus interaction. He suggested that the correction is the same as for non-interacting electrons, namely $Z^2/2$. This became one of the long standing open questions of mathematical physics (see, e.g., Lieb~\cite{Lieb1980} and Simon~\cite[Problem~10B]{Simon1984}) and was eventually proven by Siedentop and Weikard~\cite{SiedentopWeikard1986,SiedentopWeikard1987O,SiedentopWeikard1987U,SiedentopWeikard1988,SiedentopWeikard1989} (upper and lower bound) and Hughes~\cite{Hughes1986,Hughes1990} (lower bound) and later extended in various ways.

In Scott's spirit Lieb~\cite{Lieb1981} conjectured that also the density on the scale $Z^{-1}$ is given by the density of the Bohr atom. This and refinements thereof were proven by Iantchenko et al. \cite{Iantchenko1997,Iantchenkoetal1996,IantchenkoSiedentop2001}. Recently Ivrii~\cite{Ivrii2019} outlined an extension.

All these results, although mathematically correct, suffer from a serious defect viewed from a physical perspective: in the limit of large atomic numbers $Z$ the innermost electrons are attracted more and more to the nucleus. The ground state energy of such an electron is even in non-relativistic quantum mechanics already $-Z^2/2$. By the virial theorem the kinetic energy of the electron is $Z^2/2$. This means that the corresponding classical velocity is $Z$ in atomic units. This compares to the velocity of light $c$ which is $137$, a dimensionless constant. Thus, say for uranium, $Z=92$, the velocity of the innermost electrons is a substantial fraction of the velocity of light. In other words, the limit of large $Z$ renders a non-relativistic treatment questionable. A relativistic treatment is required. Comparing the energies of those electrons substantiates this view as well: the absolute value of the binding energy of the innermost electron of uranium is $4232\ Ha$ nonrelativistically compared with $8074\ Ha$ for the Dirac equation, i.e., almost a doubling. Schwinger~\cite{Schwinger1980} made this intuition quantitative and predicted a lowering of the non-relativistic Scott correction.

Analogously to the non-relativistic strong Scott conjecture by Lieb, one might predict, that the density close to the nucleus, i.e., on the scale $Z^{-1}$, behaves in a relativistic model -- after suitable renormalization -- like the sum of the absolute square of the relativistic hydrogen orbitals.

To prove such statements on the ground state energy and density starting from a microscopic model faces, however, a fundamental problem. The physically recognized starting point should be quantum electrodynamics. However even the most basic mathematical objects like the underlying Hilbert space and its Hamiltonian are unknown.

But also the straightforward generalization to a multiparticle Dirac operator -- replacing the Laplacian acting on the $n^{\rm th}$ particle by a free Dirac operator -- leads to unphysical predictions. Even if the Hamiltonian might be extended to a self-adjoint operator as recently shown by Oelker~\cite{Oelker2019} for two electrons, it leads to a spectrum which is the whole real line and dissolution of bound states, a fact that Brown and Ravenhall~\cite{BrownRavenhall1951} observed and is known as Brown--Ravenhall disease or continuum dissolution (Sucher~\cite{Sucher1980}). (See also Pilkuhn~\cite[Section~3.7]{Pilkuhn2005} for a review.)

Faced with this difficulty, various models were developed ranging from straightforward quantization of the classical relativistic Hamilton function -- which can be traced back to Chandrasekhar~\cite{Chandrasekhar1931} -- to Hamiltonians derived by physical arguments from quantum electrodynamics like the so-called no-pair Hamiltonians. All of those models have a critical coupling $\gamma=\alpha Z$ at which the energy changes from being bounded to unbounded from below (with $\alpha:=1/c$, the Sommerfeld fine structure constant). For subcritical coupling constant the Friedrichs extension yields a natural self-adjoint realization of the operator. All of them show also the above mentioned lowering of the energy.

The simplest of those models, the Chandrasekhar operator, is relatively well studied mathematically. In fact a formula for the lowering of the Scott term was proven by Solovej et al. \cite{Solovejetal2008} and Frank et al. \cite{Franketal2008}. Moreover, recently the strong Scott conjecture for the Chandrasekhar operator was proven as well (Frank et al~\cite{Franketal2020P}). However, it is known that the Chandrasekhar operator yields energies that are much too low. In fact the really heavy elements like uranium cannot by treated at the physical value of the fine structure constant, since $\alpha Z$ exceeds already $2/\pi$, the critical Chandrasekhar coupling constant.

The situation is improved for no-pair operators. (See Sucher~\cite{Sucher1980,Sucher1984,Sucher1987}; for a textbook discussion see Pilkuhn~\cite{Pilkuhn2005}.) Already the simplest, the Brown--Ravenhall operator, also called no-pair operator in the free picture, raises the energy and the critical coupling constant $2/(2/\pi+\pi/2)$ covers all known elements at the physical value of the fine structure constant. A corresponding formula for the Scott correction was obtained in~\cite{Franketal2009}. Nevertheless, its energies are still too low. \linebreak-- A convergence result for the density on the scale $Z^{-1}$ is not known.

Chemical accuracy is obtained when the external field is included in the definition of the state space. The corresponding operator is called the no-pair operator in the Furry picture. A formula for the Scott correction was proven by Handrek and Siedentop~\cite{HandrekSiedentop2015}. (The same formula should be also true when the mean field in the sense of Mittleman~\cite{Mittleman1981} is taken into account. This, however, is so far only known in Hartree--Fock approximation when the involved projection is given by the Dirac--Fock operator (Fournais et al. \cite{Fournaisetal2020}).) A formula for the ground state density on the scale $Z^{-1}$, however, is still missing. It is the purpose of this paper to close this gap and to prove the strong Scott conjecture for the no-pair operator in the Furry picture.

\section{Definitions \& main results}\label{s:defresults}
We begin with some preparatory notations which will allow to define the no-pair Hamiltonian in the Furry picture of atoms with nuclear charge $Z$ and $N$ electrons. We will use atomic units throughout, i.e., the rationalized Planck constant, the elementary charge, and the mass of the electron are all one. The energy will depend, though, besides $Z$ and $N$, also on the velocity of light $c$. For our purposes it is convenient to introduce $\gamma:=\alpha Z=Z/c$.

We write $\bp:=(1/\ri)\nabla$ for the momentum operator and
\begin{align*}
\balpha=\left(
\begin{matrix}
0 & \bsigma\\
\bsigma & 0
\end{matrix}\right),\qquad
\beta=\left(
\begin{matrix}
\one_{\cz^2} & 0\\
0 & -\one_{\cz^2}
\end{matrix}\right),
\end{align*}
with $\bsigma=(\sigma_1,\sigma_2,\sigma_3)$ the three Pauli matrices in standard representation.

We write
\begin{equation}\label{eq:defcoulombdirac}
D_{c,\,Z}:= c\balpha\cdot \bp+c^2\beta-\frac{Z}{|x|} \quad \text{in}\ L^2\left(\rz^3:\cz^4\right)
\end{equation}
for the one-electron Dirac operator defined in the sense of Nenciu~\cite{Nenciu1976} (see also \cite{KlausWust1978,Schmincke1972,Wust1975}), i.e., with form domain $H^{1/2}(\rz^3:\cz^4)$ assuming $\gamma=Z/c\in(-1,1)$.

Note
\begin{equation}\label{eq:cdrescaled}
D_{c,\,Z}\cong c^2D_{1,\,Z/c}
\end{equation}
under the scaling $x\to x/c$. For the latter we drop the first index and introduce the abbreviation
\begin{equation}\label{eq:dgamma}
D_\gamma:=D_{1,\,\gamma}.
\end{equation}
For more general electric potentials $\varphi$ allowing also for Nenciu's method to define the Dirac operator, we write
\begin{equation}\label{eq:dphi}
D_\gamma(\varphi):=D_\gamma-\varphi.
\end{equation}

Brown and Ravenhall's basic suggestion~\cite{BrownRavenhall1951} was to implement Dirac's idea \cite{Dirac1930A} of a filled Dirac sea which is inaccessible to physical electrons by requiring that the state space of an electron is the positive spectral subspace of a suitably chosen Dirac operator; in fact they suggest the positive spectral subspace of the free Dirac operator $D_{c,\,0}$. Later other choices were suggested (see Sucher~\cite{Sucher1980} for more details). A particular interesting choice yields the so called Furry picture, where the Dirac operator defining the state space is $D_{c,\,Z}$ in the atomic case. It is known that the Furry picture produces numerical values of chemical accuracy. (This choice is named after Furry, who with Oppenheimer~\cite{FurryOppenheimer1934} already introduced the corresponding splitting of the electron and positron space in quantum electrodynamics.)

In this paper we will focus on the Furry picture. To be explicit, the underlying Hilbert space is
\[
\gH_{c,\,Z}:=\underbrace{\left[\chi_{(0,\,\infty)}(D_{c,\,Z})\right]}_{=:\,\Lambda_{c,\,Z}}\left(L^2\left(\rz^3:\cz^4\right)\right).
\]
By Nenciu's above result,
\[
\Lambda_{c,\,Z}\left(\cS\left(\rz^3:\cz^4\right)\right)\subseteq H^{1/2}\left(\rz^3:\cz^4\right)
\]
and dense in $\gH_{c,\,Z}$.

The energy $\cE_{c,\,Z,\,N}$ of an atom with nuclear charge $Z$ and $N$ electrons in the state $\psi\in \bigwedge_{\nu=1}^N \Lambda_{c,\,Z}\cS(\rz^3:\cz^4)$ is
\begin{equation}\label{eq:manyfurry}
\cE_{c,\,Z,\,N}[\psi]:= \left(\psi,\left(\sum_{\nu=1}^N(\left(D_{c,\,Z}-c^2\right)_\nu+\sum_{1\,\leq\,\nu\,<\,\mu\,\leq\,N}\frac{1}{|x_\nu-x_\mu|}\right)\psi\right)
\end{equation}
The many-particle quadratic form $\cE_{c,\,Z,\,N}$ is defined whenever the one-particle quadratic form, i.e., $\cE_{c,\,Z,\,1}$ (with the understanding $\sum_{1\,\leq\,\nu\,<\,\mu\,\leq\,1}|x_\nu-x_\mu|^{-1}=0$), is defined. This is certainly true, although not necessary, if $\gamma=Z/c\in(0,1)$, an assumption which we will make throughout the rest of the paper. By construction, $\cE_{c,\,Z,\,N}$ is bounded from below and therefore extends according to Friedrichs to a closed quadratic form in the Hilbert space $\bigwedge_{n=1}^N \gH_{c,\,Z}$ with form domain $\gQ_{c,\,Z,\,N}:=\bigwedge_{n=1}^N(\Lambda_{c,\,Z}(H^{\frac12}(\rz^3:\cz^4)))$. The resulting self-adjoint operator constructed according to Friedrichs is the Furry operator of the -- possibly ionized -- atom of atomic number $Z$ with $N$ electrons.

We write $F_{c,\,Z}$ for the operator when $N=1$ and we abbreviate $\cE_{c,\,Z}:=\cE_{c,\,Z,\,Z}$ and $\gQ_{c,\,Z}:= \gQ_{c,\,Z,\,Z}$, i.e., we drop the third index of the functional and its domain, when $N=Z$.

In the one-particle case, it is also here sometimes convenient to scale out the velocity of light $c$ like for the Dirac operator and drop a factor $c^2$ with the energies. The resulting operator depends -- like in the case of the Coulomb--Dirac operator -- only on the quotient $\gamma=Z/c$ and is $F_{1,\,\gamma}$, i.e., the index pair $c,Z$ is replaced by $1,\gamma$. In this case, we simply drop the index $1$ and write $F_\gamma$ in analogy to $D_\gamma$ and similarly $\Lambda_\gamma=\Lambda_{1,\,\gamma}$. Similarly to~\eqref{eq:dphi} we introduce
\begin{equation}\label{eq:fphi}
F_\gamma(\varphi)\ \text{and}\ F_{c,\,Z}(\varphi)
\end{equation}
as the self-adjoint operators associated with $(f,(D_\gamma-\varphi-1)f)$ on $\Lambda_{\gamma}\cS(\rz^3:\cz^4)$ and $(f,(D_{c,\,Z}-\varphi-1)f)$ on $\Lambda_{c,\,Z}\cS(\rz^3:\cz^4)$ whenever closable and bounded from below.

Matte and Stockmeyer~\cite[Theorem~2.2]{MatteStockmeyer2010} showed that
\begin{equation}\label{gs}
E_{c,\,Z}:=\inf\big\{\cE_{c,\,Z}[\psi]\big| \psi\in\gQ_{c,\,Z},\; \|\psi\|=1\big\}
\end{equation}
is attained, i.e., a ground state -- not necessarily uniquely determined -- exists. Although, we neither need the state to be pure nor exactly a minimizer, we will refrain from such generalizations, and pick the state which occurs according to Lüders~\cite{Luders1951} when measuring the ground state energy, namely
\[
\Lambda:=\left(\left|\psi_1\right\rangle\left\langle\psi_1|+\dots+|\psi_D\right\rangle\left\langle\psi_D\right|\right)/D
\]
where $\psi_1,\,\dots,\,\psi_D$ is an orthonormal basis of the ground state space of $\cE_{c,\,Z}$ (with $N=Z$). We will denote the corresponding spin-summed one-particle density by
\begin{equation}\label{gsd}
\rho(x) := \frac ND \sum_{d=1}^D\;\sum_{\sigma=1}^4\int_{\Gamma^{(N-1)}} \left|\psi_d\left(x,\sigma;y\right)\right|^2\,\dy
\end{equation}
where $y\in\Gamma^{N-1}$ with $\Gamma:=\rz^3\times\{1,2,3,4\}$ are space-spin variables. Moreover, $\rd y$ is the corresponding measure, i.e., integration in the space variable and summation in the spin variable.


A refinement is to consider the density in angular momentum channels, more accurately in channels of spin-orbit coupling, labeled by
\begin{equation}\label{zp}
\kappa\in\zp:=\gz\setminus\{0\}
\end{equation}
(see Appendix~\ref{Anhang1} for more details). It is also convenient to introduce
\begin{equation}\label{l}
j_\kappa:=|\kappa|-\frac12\quad \text{and}\quad
\ell_\kappa:= j_\kappa-\frac12\sgn(\kappa)= |\kappa|-\theta(\kappa),
\end{equation}
the quantum numbers of total and orbital angular momentum all determined by $\kappa$.

The density at a point $x\in\rz^3$ in channel $\kappa\in\zp$ of the ground state $\Lambda$ of $\cE_{c,\,Z}$ is
\begin{equation}\label{eq:furryangdensity}
\rho_{\kappa}(x):= \frac{N}{4\pi D}\sum_{d=1}^D\;\sum_{\sigma\,\in\,\{+,\,-\}}\;\sum_{m=-j_\kappa}^{j_\kappa}\int_{\Gamma^{N-1}}\rd y
\left|\sum_{\tau=1}^4\int_{\bbS^2}\rd \omega\,
\overline{\Phi_{\kappa,\,m}^\sigma(\omega,\tau)}
\psi_d(|x|\omega,\tau,y)\right|^2
\end{equation}
where $\Phi^\sigma_{\kappa,\,m}$ are spherical Dirac spinors~\eqref{fi}. (See~\eqref{a1} for the relation to $\Pi_\kappa$.)

Note that $\rho=\sum_{\kappa\,\in\,\zp}\rho_\kappa$ for the state $\Lambda$. For general states, the left side needs an additional spherical average. The functions $\rho_\kappa$ and $\rho$ are the objects of interest of this work. When appropriately rescaled, we will study their convergence as $Z\to\infty$ and $Z/c$ is fixed. The objects which will turn out to be the limits are introduced now.

We write $\psi_{n,\,\kappa,\,m}$ for the orthonormal eigenfunctions of $D_\gamma$, i.e.,
\begin{equation}\label{eigenfunktion}
D_\gamma\psi_{n,\,\kappa,\,m}=\lambda_{n,\,\kappa}\psi_{n,\,\kappa,\,m}
\end{equation}
suppressing the dependence of $\gamma$ in both the eigenvalues and eigenfunctions. The corresponding eigenvalue problem was solved by Gordon~\cite{Gordon1928} and Darwin~\cite{Darwin1928}. (See also Bethe~\cite[Formula~(9.29)]{Bethe1933} or~\cite[Formula~(7.140)]{Thaller1992} for textbook treatments.)

The density of a Bohr atom for a given $\gamma$ in channel $\kappa$ is then defined by
\begin{equation}\label{eq:2.12}
\rho_{\kappa}^H(x) := \sum_{n\,=\,\theta(-\kappa)}^\infty\sum_{m=-j_\kappa}^{j_\kappa}
\sum_{\sigma=1}^4\left|\psi_{n,\,\kappa,\,m}(x,\sigma)\right|^2;
\end{equation}
the total hydrogenic density is
\begin{equation}\label{rhoh}
\rho^H(x) := \sum_{\kappa\,\in\,\zp}\rho_\kappa^H(x).
\end{equation}
Of course, this is only well defined, if the right sides of~\eqref{eq:2.12} and~\eqref{rhoh} converge which we will show outside the origin in Theorem~\ref{existencerhoh}. Moreover, we will study its behavior as $x\to0$ and $x\to\infty$.

Finally, since we show convergence in a weak -- although in fact in the radial variable rather strong -- sense, we need to specify the test functions. The test functions can be written as $U=U_1+U_2$ where $U_1$ may have a Coulomb singularity at $r=0$ and $U_2$ decays sufficiently fast as $r\to\infty$. More precisely, $U_2$ is going to belong to the test function spaces $\cD_\gamma^{(0)}$ and $\cD$ used by Frank et al. \cite{Franketal2020P}. In particular, our results imply convergence of the densities in distributional sense. For the convenience of the reader we give their definition also in Appendix~\ref{s:testfunctions}, in particular~\eqref{eq:deftestnod0} and~\eqref{eq:deftest0}. As an example, we mention that, if the test function obeys
\begin{equation}\label{eq:example}
|U(r)| \leq \const \left(r^{-1}\one_{\left\{r\,\leq\,1\right\}} + r^{-\alpha}\one_{\left\{r\,>\,1\right\}}\right),
\end{equation}
then $U=U_1+U_2$ with $U_1\in r^{-1}L^\infty_0 ([0,\infty))$ bounded and compactly supported and $U_2$ belongs to $\cD_\gamma^{(0)}$, if $\alpha>1$. It is in $\cD_\gamma^{(0)}\cap\cD$, if $\alpha>3/2$ in~\eqref{eq:example}. (The index $0$ denotes, as usual, compact support.)


\begin{theo}[Convergence of the angular momentum decomposed density]\label{convfixedl}
Fix $\kappa\in\zp$, and $U=U_1+U_2$ with $U_1\in r^{-1}L^\infty_0 ([0,\infty))$ and $U_2\in\cD_\gamma^{(0)}$. Then, with $\gamma=Z/c\in(0,1)$ fixed,
\[
\lim_{Z\,\to\,\infty}\int_{\rz^3} c^{-3}\rho_{\kappa}\left(c^{-1}x\right)U(|x|)\,\rd x = \int_{\rz^3}\rho_{\kappa}^H(x)U(|x|)\,\rd x.
\]
\end{theo}

\begin{theo}[Convergence of the density]\label{convalll}
Let $U=U_1+U_2$ with $U_1\in r^{-1}L^\infty_0 ([0,\infty))$, $U_2\in\cD\cap\cD_\gamma^{(0)}$, and $\gamma\in(0,1)$. Then
\[
\lim_{Z\,\to\,\infty}\int_{\rz^3} c^{-3}\rho\left(c^{-1}x\right)U(|x|)\,\dx = \int_{\rz^3} \rho^H(x)U(|x|)\,\dx.
\]
\end{theo}


The next result ensures that the above convergence are not meaningless. More precisely, we will now show that the hydrogenic densities are finite for all $r\in\rz_+$. To this end define for $\gamma\in[0,1]$
\begin{equation}\label{eq:defsigmagamma}
\sigma_\gamma = 1-\sqrt{1-\gamma^2}\in[0,1].
\end{equation}
Note that $\sigma_0=0$, $\sigma_1=1$, $\sigma_{\sqrt3/2}=1/2$, $\sigma_{\sqrt{15}/4}=3/4$, and $\sigma_\gamma$ is strictly monotone increasing. We will denote positive constants from now on by $A$ or $a$. A dependence on some parameter is going to be denoted by a corresponding subscript. Moreover, positive constants may vary from line to line but are still going to be denoted by the same letter.

\begin{theo}[Existence of $\rho_\kappa^H$ and $\rho^H$]\label{existencerhoh}
Let $1/2<s\leq3/4$, if $\gamma\in(0,\sqrt{15}/4)$ and $1/2<s<3/2-\sigma_\gamma$, if $\gamma\in[\sqrt{15}/4,1)$. Then there is a constant $A_{s,\,\gamma}>0$ such that for all $\kappa\in\zp$ and $x\in\rz^3\setminus\{0\}$
\begin{multline*}
\rho_{\kappa}^H(x)
\leq A_{s,\,\gamma} \frac{|\kappa|^{1-4s}}{|x|^2}\\
\times\left[\left(\frac{|x|}{|\kappa|}\right)^{2s-1}\one_{\left\{|x|\,\leq\,|\kappa|\right\}}+\left(\frac{|x|}{|\kappa|}\right)^{4s-1}\one_{\left\{|\kappa|\,\leq\,|x|\,\leq\,|\kappa|^2\right\}} + |\kappa|^{4s-1}\one_{\left\{|x|\,\geq\,|\kappa|^2\right\}}\right].
\end{multline*}
Moreover, for any $\varepsilon>0$ and $x\in\rz^3\setminus\{0\}$, there are constants $A_{\gamma,\,\varepsilon},A_\gamma>0$ such that
\begin{align}\label{eq:boundrhoh}
\rho^H(x)\leq
\begin{cases}
A_\gamma |x|^{-3/2} & \text{if}\, \gamma\in\left(0,\sqrt{15}/4\right]\\
A_{\gamma,\,\varepsilon}\left(|x|^{-2\sigma_\gamma-\varepsilon}\one_{\left\{|x|\,\leq\,1\right\}} + |x|^{-3/2}\one_{\left\{|x|\,>\,1\right\}}\right) & \text{if}\ \gamma\in\left(\sqrt{15}/4,1\right)
\end{cases}.
\end{align}
\end{theo}

Some remarks on the above results are in order.

\begin{enumerate}
\item The corresponding convergence results and pointwise bounds on the hydrogenic densities were recently proven for Chandrasekhar atoms by Frank et al. \cite{Franketal2020P}. The classes of admissible test functions are the same in both models, i.e., the test functions may have Coulomb singularities at the origin, but delta functions are not allowed, i.e., we were not able to prove pointwise convergence of the densities.

For a comparison between the results of Iantchenko et al. \cite{Iantchenkoetal1996} in the non-relativistic case with those that were obtained for the above two relativistic models, we refer to the discussion after~\cite[Theorem~1.2]{Franketal2020P}.


\item As in~\cite{Franketal2020P} we show that the hydrogenic density is finite for all $x\in\rz^3$ and obtain a pointwise upper bound with a similar asymptotic behavior for small and large distances to the nucleus. Although we are lacking a corresponding lower bound and the constant appearing in Theorem~\ref{existencerhoh} is not explicit and presumably far from sharp, we believe that the dependence on $r$ is optimal: on the one hand, relativistic effects should play a minor role for $r\gg1$ which is reflected in the $r^{-3/2}$-decay of $\rho^H$. In fact, Heilmann and Lieb~\cite{HeilmannLieb1995} proved in the non-relativistic case that the density decays like $(2^{3/2}/(3\pi^2))\gamma^{3/2}|x|^{-3/2}+o(|x|^{-3/2})$ as $x\to\infty$. Recalling that the Thomas--Fermi density satisfies $\rho_Z^{\TF}(x)=Z^2\rho_1^{\TF}(Z^{1/3}x)\sim(Z/|x|)^{3/2}$ as $x\to0$, the bounds on $\rho^H(x)$ for large $|x|$ indicate that there is a smooth transition between the quantum length scale $Z^{-1}$ and the Thomas--Fermi length scale $Z^{-1/3}$. Note also that a lower bound of the form $\rho^H(x)\geq A_\gamma |x|^{-3/2}\one_{\{|x|\,\geq\,1\}}$ would suggest that the function space $\cD_\gamma^{(0)}\cap\cD$ is optimal in the sense that it covers functions that decay like $|x|^{-3/2-\varepsilon}$, see~\eqref{eq:example}.

\item On the other hand, the behavior for small $r$ seems best possible for $\gamma\geq\sqrt{15}/4$, except for the lack of a corresponding lower bound and the arbitrary small $\varepsilon$ appearing in~\eqref{eq:boundrhoh}. The main reason for this belief is the behavior of the radial part of the hydrogenic ground state wave function at the origin
\[
\left|\psi_{n=0,\,\kappa=\,\pm1,\,m}(x)\right| \sim |x|^{\sqrt{1-\gamma^2}-1} = |x|^{-\sigma_\gamma}, \; \ m=-j_\kappa,\,\dots,\,j_\kappa.
\]
The formula reveals in particular, that the singularity of the hydrogenic density is only generated by the eigenfunctions with $|\kappa|=1$, since $\psi_{0,\,\kappa,\,m}$ has no singularity at the origin for any $|\kappa|\geq2$. This observation supports our claim for the small $r$ behavior of $\rho^H$ for $\gamma\geq\sqrt{15}/4$. However, the formula also shows that our bound for $\gamma<\sqrt{15}/4$ cannot be optimal, since it does not depend on $\gamma$ at all. As in the Chandrasekhar case, this limitation is of technical nature and comes from the restriction $\sigma_\gamma\leq3/4=\sigma_{\sqrt{15}/4}$. Ultimately, the behavior of the eigenfunctions with $|\kappa|=1$ and $\gamma=1$ suggests that the admissible singularities of our test functions are optimal. This is also expected in view of Kato's inequality since singularities which are more severe than Coulomb cannot be controlled by kinetic energy anymore.


Although the eigenfunctions $\psi_{n,\,\kappa,\,m}$ are explicitly known, the explicit summation of their absolute squares analogously to Heilmann and Lieb~\cite{HeilmannLieb1995} in the non-relativistic case is an open question. An answer would most likely allow for a more detailed study of the properties of $\rho^H$.

\item The basic idea behind the proof of the convergence result is a linear response argument which was already used by Baumgartner~\cite{Baumgartner1976}, Lieb and Simon~\cite{LiebSimon1977}, Iantchenko et al. \cite{Iantchenkoetal1996}, and Frank et al. \cite{Franketal2020R,Franketal2020P}. We first estimate the difference of the expectation values of the appropriately perturbed and unperturbed many-body Hamiltonians in the unperturbed ground state by the spectral shift between the correspondingly perturbed and unperturbed hydrogenic one-particle operators. Then, we use the generalized Feynman--Hellmann theorems~\cite[Theorem~3.2, Proposition~3.3]{Franketal2020P} to differentiate the sum of the negative eigenvalues of the perturbed hydrogenic operator. The main difficulty consists in verifying the assumptions of these theorems. In particular, we will show that the test function $U$ satisfies a certain ``relative trace class condition'' with respect to $F_\gamma$ in channel $\kappa$. To be definite let $\Pi_\kappa$ and $\Pi_\kappa^+=\Lambda_\gamma\Pi_\kappa$ be the projections onto the spaces $\ghk$ and $\ghpk=\Lambda_\gamma\ghk$ of functions with spin-orbit coupling $\kappa$ (see~\eqref{a3}-\eqref{a1} in Appendix~\ref{a:partialwaves}) and introduce the notation
\begin{equation}\label{ps}
\tr_\kappa A:= \tr\left(\Pi_\kappa A\right)
\end{equation}
when $\Pi_\kappa A$ is trace class.


For convenience we also introduce the abbreviation for the Furry operators in angular momentum channel $\kappa$
\begin{equation}\label{radialeops}
f_{\gamma,\,\kappa}:=F_\gamma|_\ghpk,\ f_{\gamma,\,\kappa}(\varphi):= F_\gamma(\varphi)|_\ghpk
\end{equation}
which we will freely use here and later.

Then our claim is that for $s>1/2$ and $\kappa\in\zp$
\[
\tr \left(f_{\gamma,\,\kappa}+1\right)^{-s}\Pi_\kappa^+U\Pi^+_\kappa\left(f_{\gamma,\,\kappa}+1\right)^{-s}<\infty.
\]
As in the Chandrasekhar case, $s>1/2$ is crucial, since $(1+k)^{-1}\notin L^1(\rz_+,\rd k)$.

The general strategy to prove the above and similar assertions, is to roll them back to those involving Chandrasekhar operators where they are known to hold~\cite[Corollary~4.5]{Franketal2020P}. A main new technical contribution is to show that the Chandrasekhar and the Furry operators are comparable: Corollary~\ref{hardydomcor} shows that one can compare $\Lambda_\gamma(|p|+1)^{2s}\Lambda_\gamma$ with $(F_\gamma+1)^{2s}$ which will be an important tool.
\end{enumerate}

\section{Applying the Feynman--Hellmann theorem in the Furry picture: the case of fixed \texorpdfstring{$\kappa$}{kappa}}

We will use the abstract version of the Hellmann--Feynman theorem by Frank et al. \cite[Theorem~3.2 and Proposition~3.3]{Franketal2020P}. To be self-contained we recall these results here. The first one will be used to handle the Coulomb singularity, whereas the second one will handle the local singularities of the test potential.

We write $A_-=-A\chi_{(-\infty,\,0)}(A)$ and denote by $\gS^1$ the set of trace class operators and by $\gS^2$ the set of Hilbert--Schmidt operators.

\begin{prop}\label{diffgen0}
Assume that $A$ is a self-adjoint operator on a Hilbert space $\gH$ with $A_-\in\gS^1(\gH)$ and $B$ is non-negative and relatively form bounded with respect to $A$. Furthermore, assume that there is $1/2\leq s\leq 1$ and $M>-\inf\sigma(A)$ such that
\begin{equation}\label{eq:traceclassdelta}
(A+M)^{-s} B (A+M)^{-s}\in\gS^1(\gH)
\end{equation}
and
\begin{equation}\label{eq:relbounddelta}
\limsup_{\lambda\,\to\,0} \left\| (A+M)^{s} (A-\lambda B+M)^{-s} \right\| <\infty.
\end{equation}
Then the one-sided derivatives $D^\pm S$ of
\[
\lambda\mapsto S(\lambda):=\tr(A-\lambda B)_-
\]
satisfy
\[
\tr B\chi_{(-\infty,\,0)}(A)=D^-S(0)\leq D^+S(0)=\tr B\chi_{(-\infty,\,0]}(A).
\]
In particular, $S$ is differentiable at $\lambda=0$, if and only if $B|_{\ker A}=0$.
\end{prop}

\begin{prop}\label{diffgen}
Assume that $A$ is self-adjoint with $A_{-} \in\gS^1$ and $B$ is non-negative, and let $1/2<s\leq1$. Assume that there is an $s'<s$ such that for some $M>-\inf\sigma(A)$, \eqref{eq:traceclassdelta} holds, and that for some $a>0$
\begin{equation}\label{eq:increasedrelbdd}
B^{2s} \leq a(A+M)^{2s'}.
\end{equation}
Then $B$ is form bounded with respect to $A$ with form bound $0$ and the conclusions in Proposition~\ref{diffgen0} hold.
\end{prop}

We recall the following two observations.
\begin{enumerate}
\item By~\eqref{eq:traceclassdelta},
\[
\tr\left(B\chi_{(-\infty,\,0]}(A)\right)=-\tr\left[\big((A+M)^{-s}B(A+M)^{-s}\big)\left(\chi_{(-\infty,\,0]}(A)(A+M)^{2s}\right)\right]<\infty,
\]
since $\chi_{(-\infty,\,0]}(A)(A+M)^{2s}$ is bounded. Hence, also $D^+S(0)<\infty$.

\item If the bottom of the essential spectrum of $A$ is strictly positive, the result coincides with the classical Feynman--Hellmann theorem. The point is that the formulae remain valid even, if $\inf\sigma_{\ess}(A)=0$, i.e., the case where perturbation theory is not directly applicable.
\end{enumerate}
In the application of the two propositions above, the underlying Hilbert space is $\ghpk$, $A$ will be the Furry operator $F_{\gamma}$ restricted to this space, and $B = \Lambda_\gamma(U\otimes\one_{\cz^4})\Lambda_\gamma$ also restricted to this space plays the role of the test potential.

We recall some basic facts about $F_\gamma$.
\begin{lemm}\label{hydrogen}
Let $\kappa\in\zp$ and $\gamma\in[0,1)$. Then $F_\gamma\geq\sqrt{1-\gamma^2}-1$, $0\notin\sigma_\pp(F_\gamma)$, and $F_\gamma+1$ has a bounded inverse. Moreover, $\tr_\kappa (F_\gamma)_-<\infty$.
\end{lemm}


\begin{proof}
The fact that $\rz\setminus(-1,1)\cap \sigma_\pp(D_\gamma)=\emptyset$ is a standard consequence of the virial theorem proven by Kalf~\cite{Kalf1976} for all $\gamma\in(-1,1)$. In particular the eigenvalues of $D_\gamma$ are all given by Sommerfeld's eigenvalue formula
\begin{equation}\label{eq:eigenvalue}
\lambda_{n,\,\kappa} = \left(1+\frac{\gamma^2}{\left(n+\sqrt{\kappa^2-\gamma^2}\right)^2}\right)^{-1/2}
\end{equation}
with $(\kappa,n)\in(-\nz\times\nz)\cup(\nz\times\nz_0)$ (Sommerfeld~\cite{Sommerfeld1916}, Gordon~\cite{Gordon1928}, and Darwin~\cite{Darwin1928}). In particular the lowest eigenvalue is $\sqrt{1-\gamma^2}$ and $\sum_{n}(\lambda_{n,\,\kappa}-1)$ is absolutely summable for each fixed $\kappa$.
\end{proof}

For a textbook discussion of $D_\gamma$, we refer to Bethe~\cite{Bethe1933} and Thaller~\cite{Thaller1992}, in particular~\cite[p.~314f]{Bethe1933} and~\cite[Sections~7.4.2 and~7.4.5]{Thaller1992} for the discussion of the point spectrum.

The following two propositions show the applicability of Propositions~\ref{diffgen0} and~\ref{diffgen}. They are the keys to prove Theorem~\ref{convfixedl} and will be proved in Subsections~\ref{ss:traceinequalities} and~\ref{ss:controlingcoulomb}.

\begin{prop}\label{diffu2}
Let $\gamma\in(0,1)$, $\kappa\in\zp$, and $0\leq U\in\cD_\gamma^{(0)}$. Then
\[
\lambda\mapsto \tr f_{\gamma,\,\kappa}(\lambda U)_-
\]
is differentiable at $\lambda=0$ with derivative $\int_{\rz^3}\rho_\kappa^H(x) U(|x|)\,\dx$.
\end{prop}

\begin{prop}\label{diffu1}
Let $\gamma\in(0,1)$, $\kappa\in\zp$, and $0\leq U \in r^{-1} L^\infty_0 ([0,\,\infty))$. Then
\[
\lambda\mapsto
\tr f_{\gamma,\,\kappa}(\lambda U)_-
\]
is differentiable at $\lambda=0$ with derivative $\int_{\rz^3} \rho_\kappa^H(x) U(|x|)\,\dx$.
\end{prop}

Note that Propositions~\ref{diffu2} and~\ref{diffu1} imply $\int_{\rz^3}\rho_\kappa^H(x)U(|x|)\dx<\infty$. In particular, these results show that for $\gamma<1$ and $R>0$,
\[
\int_{|x|\,<\,R} \rho_\kappa^H(x) |x|^{-1}\,\dx < \infty.
\]
In fact, there is also a simple, direct proof of this, even when $R=\infty$: based on a computation of Burke and Grant~\cite{BurkeGrant1967}, Handrek and Siedentop~\cite[Lemma~2]{HandrekSiedentop2015} show that the potential energy of hydrogenic eigenfunctions satisfies
\[
\left(\psi_{n,\,\kappa,\,m},\frac{\gamma}{|x|}\psi_{n,\,\kappa,\,m}\right)
\leq\frac{a_{\gamma_0}\gamma^2}{(n+|\kappa|)^2}.
\]
Clearly, the right side is summable in $n$ and -- trivially -- in $m$.

Propositions~\ref{diffu2} and~\ref{diffu1} will be deduced from Propositions~\ref{diffgen} and~\ref{diffgen0} respectively. To verify their assumptions, we will first reduce the problem to the scalar Chandrasekhar operator and then use~\cite{Franketal2020P}.

In this and the next section, we often use the Davis--Sherman inequality (Davis \cite{Davis1957}) and a closely related version thereof by Frank and Geisinger \cite[Lem\-ma~6.4]{FrankGeisinger2016}.

\begin{lemm}
Let $P$ be an orthogonal projection, $A\geq0$ be a linear operator, $s>0$, and $[0,\infty)\ni x\mapsto f(x):=x^s$. Then the following statements hold:
\begin{enumerate}
\item If $s\in[1,2]$, then one has the quadratic form inequality
\begin{equation}\label{eq:SD}
f(PAP) = Pf(PAP)P \leq Pf(A)P\,.
\end{equation}

\item If $s\in(0,1]$ and the kernel of $A$ is trivial, then one has the quadratic form inequality
\begin{equation}\label{eq:FG}
f(PAP) = Pf(PAP)P \geq Pf(A)P\,.
\end{equation}
\end{enumerate}
\end{lemm}

\begin{rema}
\leavevmode
\begin{enumerate}
\item The inequality $Pf(PAP)P \geq Pf(A)P$ in~\eqref{eq:FG} holds true for all operator monotone functions $f$. Similarly, and more generally, the inequality $Pf(PAP)P\leq Pf(A)P$ in~\eqref{eq:SD} extends to all self-adjoint operators $A$ and operator convex functions $f$ which are defined on the spectrum of $A$. Moreover, operator convexity is also necessary for the inequality $Pf(PAP)P\leq Pf(A)P$ in~\eqref{eq:SD}.

\item The equality $f(PAP)=Pf(PAP)P$ in~\eqref{eq:SD} and~\eqref{eq:FG} holds whenever, e.g., $f(0)=0$. This is a consequence of the spectral theorem, cf.~\cite[p.~43]{Davis1957}.

\item Inequality~\eqref{eq:SD} was observed by Sherman for $f(x)=x^2$ and later proved by Davis~\cite{Davis1957} for all operator convex functions. See also Bhatia~\cite[Theorem~V.2.3]{Bhatia1997} and Carlen~\cite[Theorem~4.19]{Carlen2010} for textbook references. Inequality~\eqref{eq:FG} was proved by Frank and Geisinger~\cite[Lemma~6.4]{FrankGeisinger2016}.
\end{enumerate}
\end{rema}



\subsection{Comparison between the Chandrasekhar and the Furry operator}

We write $D^0_\gamma:=\balpha\cdot\bp-\gamma|x|^{-1}$ for the massless Coulomb--Dirac operator (which is defined as in Section~\ref{s:defresults}, Nenciu~\cite{Nenciu1976}). The following lemma gives a comparison between $|p|^s$ and $|D_\gamma^0|^s$ as operators in $L^2(\rz^3:\cz^4)$.


\begin{lemm}[{Frank et al.~\cite[Corollary~1.8]{Franketal2021}}]\label{hardydom}
Let $\gamma\in[0,1)$ and $s\in(0,1]$. Then there exists an $A_{s,\,\gamma}<\infty$ such that
\[
\left|D^0_\gamma\right|^{2s} \leq A_{s,\,\gamma} |p|^{2s}\,.
\]
If, additionally, $s<3/2-\sigma_{\gamma}$, then there is an $a_{s,\,\gamma}>0$ such that
\[
\left|D^0_\gamma\right|^{2s} \geq a_{s,\,\gamma} |p|^{2s}\,.
\]
\end{lemm}

From Lemma~\ref{hardydom}, we deduce

\begin{coro}\label{hardydomcor}
Let $\gamma\in[0,1)$ and $M\geq0$. If $0<s<\min\{3/2-\sigma_\gamma,1\}$, then
\begin{equation}\label{eq:hardydomcorform}
\Lambda_\gamma \left(|p|^{2s}+M\right)\Lambda_\gamma
\leq \left(1-\gamma^2\right)^{-s}\left(a_{s,\,\gamma}^{-1}+M\right) \left(F_\gamma+1\right)^{2s}.
\end{equation}
Moreover, if $0<s\leq1$, then
\[
\left(F_\gamma+1\right)^{2s}
\leq 2^s\left(1+4\gamma^2\right)^s \Lambda_\gamma\left(|p|+\left(1+4\gamma^2\right)^{-\frac12} \right)^{2s}\Lambda_\gamma.
\]
\end{coro}

\begin{proof}
We begin with the first claim. Since $\sqrt{1-\gamma^2}$ is the lowest positive spectral point of $D_\gamma$, it suffices to show the claim for $M=0$. Next, note that
\[
\left(D_\gamma\right)^2 \geq \left(1-\gamma^2\right) \left|D_\gamma^0\right|^2
\]
by Morozov and M\"uller~\cite[Proof of Corollary~I.2]{MorozovMuller2017}. By operator monotonicity of $x\mapsto x^s$ with $s\in(0,1]$, and Lemma~\ref{hardydom} we have
\begin{multline*}
\Lambda_\gamma|p|^{2s}\Lambda_\gamma \leq \Lambda_\gamma a_{s,\,\gamma}^{-1}\left|D_\gamma^0\right|^{2s}\Lambda_\gamma
\leq \left(1-\gamma^2\right)^{-s}a_{s,\,\gamma}^{-1}\Lambda_\gamma\left|D_\gamma\right|^{2s}\Lambda_\gamma\\
 = \left(1-\gamma^2\right)^{-s}a_{s,\,\gamma}^{-1}\left(\Lambda_\gamma D_\gamma\Lambda_\gamma\right)^{2s}
\end{multline*}
where the last step is obvious by the spectral theorem.

We turn to the second inequality. First we note that the left side is equal to $\Lambda_\gamma |D_\gamma|^{2s}\Lambda_\gamma$, i.e., it suffices to prove the stronger inequality
\[
\left|D_\gamma\right|^{2s}\leq 2^s\left(1+4\gamma^2\right)^s\left(|p|+\left(1+4\gamma^2\right)^{-1/2}\right)^{2s}\,.
\]
By operator monotonicity of roots, it is enough to prove the claim for the largest occurring $s$, namely $s=1$. This, however, follows by first using the Schwarz inequality and then Hardy's inequality
\begin{align*}
\left|D_\gamma\right|^2 \leq 2\left(p^2+1+\frac{\gamma^2}{|x|^2}\right)
&\leq 2 \left(|p|^2\left(1+4\gamma^2\right)+1\right)\\ &\leq 2\left(1+4\gamma^2\right)\left(|p|+\frac{1}{\sqrt{1+4\gamma^2}}\right)^2
\end{align*}
where the last step is obvious.
\end{proof}

We introduce the following restricted operators in $\gh_{\kappa}$,
\begin{align*}
p_{\ell_\kappa} & := \sqrt{-\Delta}|_{\gh_{\kappa}},\\
C_{\ell_\kappa} & := \left(\sqrt{-\Delta+1}-1\right)\big|_{\gh_{\kappa}}.
\end{align*}
The corresponding radial operators in $L^2(\rz_+:\cz,\dr)$ are
\begin{align}\label{eq:radialkinetic}
p_{\ell_\kappa}^{(r)} & := \sqrt{ - \frac{\rd^2}{\dr^2} + \frac{\ell_\kappa\left(\ell_\kappa+1\right)}{r^2}},\\
C_{\ell_\kappa}^{(r)} & :=\sqrt{- \frac{\rd^2}{\dr^2} + \frac{\ell_\kappa\left(\ell_\kappa+1\right)}{r^2} + 1}-1.
\end{align}
We note that the bounds of Corollary~\ref{hardydomcor} continue to hold in each $\gh_{\kappa,\,m}$. Recall that any element $f\in\gh_{\kappa,\,m}$ is of the form
\[
f(x)=\sum_{\sigma\,\in\,\{+,\,-\}} |x|^{-1} f^\sigma(|x|)\Phi_{\kappa,\,m}^\sigma(x/|x|)
\]
where the $\Phi_{\kappa,\,m}^\sigma$ are defined in~\eqref{fi} and $f^\sigma\in L^2(\rz_+)$. Both $D_\gamma$ and $|p|$ leave the spaces $\gh_{\kappa,\,m}$ invariant, i.e., they commute with the projection $\Pi_{\kappa,\,m}$. Indeed, for $f\in\gh_{\kappa,\,m}\cap H^{1}(\rz^3:\cz^4)$ and $g\in\gh_{\kappa',\,m'}\cap H^{1}(\rz^3:\cz^4)$, one has

\[
\left(f,|p| g\right)_{L^2\left(\rz^3\,:\,\cz^4\right)} = \left(\left(f^+,p_{\ell_\kappa}^{(r)} g^+\right)_{L^2\left(\rz_+\,:\,\cz\right)} + \left(f^-,p_{2j_\kappa-\ell_\kappa}^{(r)} g^-\right)_{L^2\left(\rz_+\,:\,\cz\right)}\right)\delta_{\kappa,\,\kappa'}\delta_{m,\,m'}
\]
and
\begin{multline}\label{eq:radialdirac}
(f,D_\gamma g)_{L^2\left(\rz^3\,:\,\cz^4\right)}\\
= \left(\left(
\begin{matrix}
f^+\\ f^-
\end{matrix}
\right),\left(
\begin{matrix}
1-\frac{\gamma}{r} & -\frac{\rd}{\dr}-\frac{\kappa}{r}\\
\frac{\rd}{\dr}-\frac{\kappa}{r} & -1-\frac{\gamma}{r}
\end{matrix}
\right)\left(
\begin{matrix}
g^+\\ g^-
\end{matrix}
\right)\right)_{L^2\left(\rz_+\,:\,\cz^2\right)}\delta_{\kappa,\,\kappa'}\delta_{m,\,m'},
\end{multline}
see also~\cite[Formula~(7.105)]{Thaller1992}. Together with the spectral theorem, this shows that the projection of~\eqref{eq:hardydomcorform} onto $\gh_{\kappa,\,m}$, namely
\begin{align}
\Pi_{\kappa,\,m}\Lambda_\gamma(|p|+M)^{2s}\Lambda_\gamma\Pi_{\kappa,\,m}
&\leq a_{s,\,\gamma} \Pi_{\kappa,\,m}\left(F_\gamma+1\right)^{2s}\Pi_{\kappa,\,m},\label{eq:hardydomcorang0}\\
\intertext{is equivalent to}
\Lambda_\gamma\left(\Pi_{\kappa,\,m}(|p|+M)\Pi_{\kappa,\,m}\right)^{2s}\Lambda_\gamma
&\leq a_{s,\,\gamma}\big(\Pi_{\kappa,\,m}\left(F_{\gamma}+1\right)\Pi_{\kappa,\,m}\big)^{2s}.
\label{eq:hardydomcorang}
\end{align}
Mutatis mutandis, the equivalence holds also for the projection onto $\gh_\kappa$.

\subsection{Trace inequalities in \texorpdfstring{$\gh_\kappa$}{gh kappa}}\label{ss:traceinequalities}

We recall some trace and form inequalities for functions belonging to the spaces $\cK_{s}^{(0)}$ introduced by Frank et al. \cite{Franketal2020P}. For the convenience of the reader, we give their definition in Appendix~\ref{s:testfunctions}. Frank et al. \cite{Franketal2020P} wrote the associated trace inequalities in terms of powers of $C_\ell^{(r)}+M$. Using Plancherel's theorem, one can rewrite them as inequalities in powers of $p_\ell^{(r)}+M$ instead. Here, we will actually formulate the inequalities in terms of $p_\ell$.

\begin{lemm}\label{genReltrclassnod}
Let $M>0$, $s\in(1/2,1]$, $\kappa\in\zp$, and $0\leq W \in\cK_{s}^{(0)}$. Then
\begin{equation}\label{eq:genReltrclassnod}
\tr\left[\left(p_{\ell_\kappa}+M\right)^{-s}W\left(p_{\ell_\kappa}+M\right)^{-s}\right]
\leq A_{s,\,\kappa,\,M} \|W\|_{\cK_{s}^{(0)}}.
\end{equation}
In particular, we have in $L^2(\rz^3:\cz^4)$
\begin{equation}\label{eq:genSobolevDualnod}
\Pi_{\kappa} W\Pi_{\kappa}
\leq A_{s,\,\kappa,\,M} \|W\|_{\cK_{s}^{(0)}} \big(\Pi_{\kappa}(|p|+M)\Pi_{\kappa}\big)^{2s}.
\end{equation}
\end{lemm}

\begin{proof}
The estimate~\eqref{eq:genReltrclassnod} follows from
\[
\left\|W^{\frac12}\left(p_{|\kappa|}^{(r)}+M\right)^{-s}\right\|_{\gS^2\left(L^2\left(\rz_+\right)\right)}^2 + \left\|W^{\frac12}\left(p_{|\kappa|-1}^{(r)}+M\right)^{-s}\right\|_{\gS^2\left(L^2\left(\rz_+\right)\right)}^2
\leq A_{s,\,\kappa,\,M} \|W\|_{\cK_{s}^{(0)}}
\]
(Frank et al. \cite[Proposition~4.4]{Franketal2020P}). Estimate~\eqref{eq:genSobolevDualnod} follows immediately from \eqref{eq:genReltrclassnod}.
\end{proof}

Combining Corollary~\ref{hardydomcor} in each channel $\kappa$, i.e., \eqref{eq:hardydomcorang} and Lemma~\ref{genReltrclassnod} yields a generalization of the previous inequalities but now with respect to the Furry operator. Using the notation $\Pi^+_\kappa$ defined in~\eqref{a2} we have
\begin{lemm}\label{genReltrclassHydro}
For $\gamma\in(0,1)$, $1/2<s<\min\{3/2-\sigma_\gamma,1\}$, $\kappa\in\zp$, and \linebreak$0\leq W \in\cK_{s}^{(0)}$
\begin{equation}\label{eq:genReltrclassHydro}
\tr\left[\left(f_{\gamma,\,\kappa}+1\right)^{-s}\Pi^+_\kappa W\Pi^+_{\kappa} \left(f_{\gamma,\,\kappa}+1\right)^{-s}\right]
\leq A_{\gamma,s,\,\kappa} \|W\|_{\cK_{s}^{(0)}}.
\end{equation}
In particular
\begin{equation}\label{eq:genSobolevDualHydro}
\Pi_{\kappa}^+W\Pi_{\kappa}^+
\leq A_{\gamma,\,s,\,\kappa} \|W\|_{\cK_{s}^{(0)}} \left(f_{\gamma,\,\kappa}+1\right)^{2s}.
\end{equation}
\end{lemm}


\begin{proof}[Proof of Proposition~\ref{diffu2}]
We apply Proposition~\ref{diffgen} with $M=1$, $A = f_{\gamma,\,\kappa}$, and $B = \Pi_\kappa^+ U\Pi^+_\kappa$. Here $U\in\cK_{s}^{(0)}$ and $U^{2s}\in\cK_{s'}^{(0)}$ with $1/2<s'<s\leq1$, if $\gamma<\sqrt3/2$, and $1/2<s'<s<3/2-\sigma_\gamma$, if $\gamma\in[\sqrt3/2,1)$.

We now verify the assumptions of Proposition~\ref{diffgen}: the assumptions on $F_{\gamma}$ in $\gH_\kappa^+$ follow from Lemma~\ref{hydrogen}. In particular, since zero is not an eigenvalue of $F_{\gamma}$, the right and left derivative agree at $\lambda=0$.

Since $U\in\cK_{s}^{(0)}$, we have by Lemma~\ref{genReltrclassHydro} that $(\ppp U\ppp)^{1/2}(f_{\gamma,\,\kappa}+1)^{-s}\in \gS^2(\ghpk).$


Eventually, we verify
\begin{equation}\label{eq:NZ}
\left(\Pi_\kappa^+ U \Pi_\kappa^+\right)^{2s}
\leq a_{\gamma,s,\,\kappa}\left(f_{\gamma,\,\kappa}+1\right)^{2s'}
\end{equation}
for $1/2<s'<s$. To show~\eqref{eq:NZ}, we use~\eqref{eq:SD} with $f(x)=x^{2s}$ and obtain
\[
\left(\Pi^+_\kappa U \ppp\right)^{2s}
\leq \Pi_\kappa^+ U^{2s} \Pi_\kappa^+.
\]
Hence, by~\eqref{eq:genSobolevDualHydro}, the left side of~\eqref{eq:NZ} is bounded by $(f_{\gamma,\,\kappa}+1)^{2s'}$ times a constant, since $U^{2s}\in\cK_{s'}^{(0)}$.
\end{proof}

\subsection{Controlling Coulomb perturbations}\label{ss:controlingcoulomb}

The main difficulty in applying Proposition~\ref{diffgen0} is verifying~\eqref{eq:relbounddelta}. In our setting it is an inequality for each fixed $\kappa\in \zp$. However, it follows from the following stronger statement which does not need a partial wave analysis.
\begin{lemm}\label{boundpertcoulomb}
Let $\gamma\in(0,1)$, $0\leq U\in r^{-1}L^\infty([0,\infty))$, and $\tfrac12\leq s<\min\{\tfrac32-\sigma_\gamma,1\}$. Then there is a $a_{\gamma,\,s}\in\rz$ and a $\lambda_0>0$ such that for $|\lambda|<\lambda_0$
\begin{equation}
\left(F_\gamma+1\right)^{2s}
\leq a_{\gamma,\,s} \left(F_{\gamma}(\lambda U)+1\right)^{2s}.
\end{equation}
\end{lemm}


\begin{proof}
Since $F_\gamma>-1$, obviously, for sufficiently small $\lambda$, $F_{\gamma}+1 -\lambda \Lambda_\gamma U \Lambda_\gamma>0$.

Next we first assume $\lambda>0$. By operator convexity of $x\mapsto x^{2s}$ with $s\in[1/2,1]$ and the Davis--Sherman inequality~\eqref{eq:SD}, we obtain
\begin{equation}\label{eq:convexity}
\begin{split}
\left(F_\gamma+1\right)^{2s} &= \left(F_{\gamma}(\lambda U)+1 + \lambda \Lambda_\gamma U\Lambda_\gamma \right)^{2s}\\
&\leq 2^{2s-1}\left(F_{\gamma}(\lambda U)+1 \right)^{2s}
+ 2^{2s-1}\lambda^{2s}\left(\Lambda_\gamma U\Lambda_\gamma \right)^{2s}\\
&\leq 2^{2s-1}\left(F_{\gamma}(\lambda U)+1 \right)^{2s} + 2^{2s-1}\lambda^{2s}\Lambda_\gamma U^{2s}\Lambda_\gamma.
\end{split}
\end{equation}
Since $U(r)\leq \|rU\|_\infty/r$, Hardy's inequality yields $ U^2 \leq 4\|rU\|_\infty^2|p|^2$. Thus, by operator monotonicity of roots and Corollary~\ref{hardydomcor},
\[
\Lambda_\gamma U^{2s}\Lambda_\gamma \leq 4^s\left\|rU\right\|_\infty^{2s}\Lambda_\gamma |p|^{2s}\Lambda_\gamma
\leq 4^s\left\|rU\right\|_\infty^{2s}A_{s,\,\gamma}\left(F_\gamma+1\right)^{2s}
\]
where $A_{s,\gamma}$ is the constant in~\eqref{eq:hardydomcorform}. Plugging this estimate in~\eqref{eq:convexity} yields

\[
\left(F_\gamma+1\right)^{2s} \leq 2^{2s-1}\left(1-2^{4s-1}A_{s,\,\gamma}\left\|rU\right\|_\infty^{2s}\lambda^{2s}\right)^{-1}
\left(F_{\gamma}(\lambda U)+1\right)^{2s}
\]
proving the assertion for $\lambda>0$.

If $\lambda<0$, we set $\mu:=-\lambda>0$ and $\varepsilon:=\sqrt\mu\in(0,1)$. By operator convexity
\begin{align}
\left(F_\gamma+1\right)^{2s} = &\left[(1-\varepsilon)\left(F_\gamma\left(-\mu U\right)+1\right) + \varepsilon\left(F_\gamma\big((1-\varepsilon)\varepsilon^{-1}\mu U\big)+1 \right)\right]^{2s}\label{eq:convexity2}\\
\leq &2^{2s-1}(1-\varepsilon)^{2s}\big(F_\gamma\left(-\mu U\right)+1\big)^{2s}\label{40a}\\
& + 2^{2s-1}\varepsilon^{2s}\left(F_\gamma\big((1-\varepsilon)\varepsilon^{-1}\mu U\big)+1\right)^{2s}.\label{40b}
\end{align}
Here, we used that both operators are non-negative (because of the condition on $\lambda=-\mu$ which implies that the coupling constant of the perturbation in the second summand is $\cO(\sqrt\mu)$ which is chosen sufficiently small). This allows to use operator convexity in~\eqref{40a}. Suppose there is an $\upsilon\in\rz$ and a $\mu_0>0$ such that for all $\mu\in[0,\mu_0]$
\begin{equation}\label{40c}
~\eqref{40b} \leq \varepsilon^{2s}\upsilon\left(F_\gamma+1\right)^{2s},
\end{equation}
then, the assertion follows as in the case $\lambda>0$ by taking this term to the left side of~\eqref{eq:convexity2} and dividing both sides of the inequality by $(1-\varepsilon^{2s}\upsilon)$, which is allowed for sufficiently small $\varepsilon$.

We turn to the proof of~\eqref{40c}: since $\Lambda_\gamma D_\gamma=\Lambda_\gamma|D_\gamma|$ and by~\eqref{eq:SD}, we have
\[
\left(\Lambda_\gamma\left(D_\gamma-\frac{1-\varepsilon}{\varepsilon}\mu U\right)\Lambda_\gamma\right)^{2s} \leq \Lambda_\gamma\left(\left|D_\gamma\right|-\frac{1-\varepsilon}{\varepsilon}\mu U\right)^{2s}\Lambda_\gamma.
\]
Now, $(\balpha\cdot\bp+\beta)^2=p^2+1$, $U(r)\leq\tfrac1r\|rU\|_\infty $, Schwarz' and Hardy's inequality imply
\begin{align*}
\left(\left|D_\gamma\right|-\frac{1-\varepsilon}{\varepsilon}\mu U\right)^2 & \leq 4\left(|p|^2+1+\left(\gamma+\frac{1-\varepsilon}{\varepsilon}\mu\left\|rU\right\|_\infty\right)^2|x|^{-2}\right)\\
& \leq 4\left(\left(1+4\left(\gamma+\frac{1-\varepsilon}{\varepsilon}\mu\left\|rU\right\|_\infty\right)^2\right)|p|^2+1\right)\\
& \leq 4\left(1+4\left(\gamma+\frac{1-\varepsilon}{\varepsilon}\mu\left\|rU\right\|_\infty\right)^2\right)\left(|p|+1\right)^2.
\end{align*}
Using $(|p|+1)^{2s}\leq2^{2s-1}(|p|^{2s}+1)$ for $s\in[1/2,1]$, and applying Corollary~\ref{hardydomcor}, we see the existence of the wanted $\upsilon$, if $\mu$, and thus $\mu/\epsilon$, is sufficiently close to zero.
\end{proof}


\begin{proof}[Proof of Proposition~\ref{diffu1}]
We apply Proposition~\ref{diffgen0} to the operators $ A = f_{\gamma,\,\kappa}$ and $B =\ppp U \ppp\geq 0$ with $1/2<s<\min\{3/2-\sigma_\gamma,1\}$.

We have already verified the assumptions concerning $f_{\gamma,\,\kappa}$ in the proof of Proposition~\ref{diffu2}. The fact that $\ppp U\ppp$ is relatively form bounded with respect to $f_{\gamma,\,\kappa}$ follows from Kato's inequality and Corollary~\ref{hardydomcor} in every channel $\kappa$.

Since $U\in\cK_{s}^{(0)}$, Lemma~\ref{genReltrclassHydro} implies
\[
\tr\left(f_{\gamma,\,\kappa}+1\right)^{-s}\ppp U\ppp\left(f_{\gamma,\,\kappa}+1\right)^{-s}<\infty.
\]
Finally, Assumption~\eqref{eq:relbounddelta} follows from
\[
\left(F_\gamma+1\right)^{2s}\leq A_{\gamma,\,s}\left(F_{\gamma}(\lambda U)+1\right)^{2s}
\]
for small $|\lambda|$ which is Lemma~\ref{boundpertcoulomb}.
\end{proof}



\section{Controlling large angular momenta}\label{s:alll}

\subsection{Estimating the spectral shift in channel \texorpdfstring{$\kappa$}{kappa}}
We will use the notation $D_\gamma(\varphi)$ and $F_\gamma(\varphi)$ introduced in~\eqref{eq:dphi} and~\eqref{eq:fphi}. The following proposition will allow to apply the Weierstra\ss~M-test to deduce Theorem~\ref{convalll} from Theorem~\ref{convfixedl} as in~\cite{Franketal2020P}.

\begin{prop}\label{dominantsigma}
Let $0<\gamma<1$, $0\leq V(r)\leq\gamma/r$ for $r\in\rz_+$, and $U=U_1+U_2$ with $0\leq U_1\in r^{-1}L^\infty_0 (\rz_+)$ and $0\leq U_2\in\cD$. Then there are $\varepsilon>0$, $A_{\gamma,s}<\infty$, $\lambda_1>0$, and $K_\gamma\in\nz$ such that $|\lambda|<\lambda_1$ and $|\kappa|\geq K_\gamma$ implies
\begin{equation}
\tr_\kappa F_0(V+\lambda U)_- -\tr_\kappa F_0(V)_-
\leq A_{\gamma,\,s}\lambda \|U\|_{\cK_{s,\,0}} |\kappa|^{-1-\varepsilon}.
\end{equation}
\end{prop}

In preparation of the proof, we give a trace and a Sobolev inequality with respect to $C_{\ell_\kappa}+a \kappa^{-2}$ on $\gh_\kappa$. We recall that these inequalities were crucial in~\cite{Franketal2020P} to treat functions belonging to $\cK_{s,\,\delta}$. The following lemma follows from~\cite[Proposition~5.2]{Franketal2020P} in the same way as Lemma~\ref{genReltrclassnod} followed from~\cite[Proposition~4.4]{Franketal2020P}.

\begin{lemm}\label{genReltrclass}
Let $a>0$, $\delta\in[0,2s-1]$, $s\in(1/2,3/4]$, $\kappa\in\zp$, and $0\leq W \in\cK_{s,\delta}$. Then,
\begin{equation}\label{eq:genReltrclass}
\left\|W^{1/2}\left(C_{\ell_\kappa}+a\kappa^{-2}\right)^{-s}\right\|_{\gS^2\left(\gh_{\kappa}\right)}^2
\leq A_{s,\,a} |\kappa|^{1-\delta}\|W\|_{\cK_{s,\,\delta}}.
\end{equation}
In particular, in $L^2(\rz^3:\cz^4)$,
\begin{equation}\label{eq:genSobolevDual}
\Pi_{\kappa}W\Pi_{\kappa}
\leq A_{s,\,a} |\kappa|^{-\delta}\|W\|_{\cK_{s,\,\delta}} \left(\Pi_{\kappa}\left(\sqrt{-\Delta+1}-1+a\kappa^{-2}\right)\Pi_{\kappa}\right)^{2s}.
\end{equation}
\end{lemm}

Note, that the latter inequality follows immediately, since the multiplicity of each eigenvalue is proportional to $|\kappa|$.

To prove Proposition~\ref{dominantsigma}, we will again control Dirac operators by scalar operators:
\begin{lemm}\label{boundkineticenergy}
Let $a>0$ and $\kappa\in\zp$ such that $1\geq a \kappa^{-2}$. Then
\[
\left(D_0-1+a \kappa^{-2}\right)^2 \geq \left(\sqrt{p^2+1}-1+a \kappa^{-2}\right)^2.
\]
\end{lemm}

\begin{proof}
The assertion is equivalent to the inequality
\[
p^2+1+\left(1-\frac a{\kappa^2}\right)^2-2\left(1-\frac a{\kappa^2}\right)D_0\\
\geq p^2+1+\left(1-\frac a{\kappa^2}\right)^2-2\left(1-\frac a{\kappa^2}\right)\sqrt{p^2+1}.
\]
Since $1\geq a \kappa^{-2}$ and $D_0=\balpha\cdot\bp+\beta \leq |\balpha\cdot\bp+\beta|=\sqrt{p^2+1}$, the assertion follows.
\end{proof}

\begin{lemm}\label{boundA}
Let $\gamma\in(0,1)$, $0\leq V(r)\leq\gamma/r$, $s\in(1/2,3/4]$, and $U=U_1+U_2$ with $U_1\in r^{-1}L^\infty([0,\infty))$ and $|U_2|^{2s}\in \cK_{s,\,0}$. Then there are constants $K_{\gamma}\in\nz$ and $a_\gamma,\lambda_2\in \rz_+$, such that for all $\lambda\in[-\lambda_2, \lambda_2]$ and all $\kappa\in\zp$ with $|\kappa|\geq K_\gamma$
\begin{equation}\label{agamma}
f_{0,\,\kappa}(V+\lambda U) \geq -a_\gamma \kappa^{-2}
\end{equation}
holds.
\end{lemm}


\begin{proof}
Note that Sommerfeld's eigenvalue formula~\eqref{eq:eigenvalue} immediately implies \eqref{agamma} for pure Coulomb potentials $V(r)+\lambda U(r)=\gamma/r$. In this case it will be useful to emphasize the Coulombic origin and write $a_{\gamma,\,C}$ instead of $a_\gamma$.

Since
\[
\int_0^\infty\left(\ell+\frac12\right)^2r^{-2}|g(r)|^2\,\dr \leq \left(g,\left(p_\ell^{(r)}\right)^2g\right)_{L^2\left(\rz_+\right)}
\]
(Hardy) and by picking $\ell=\ell_\kappa=|\kappa|-\theta(\kappa)$ (see~\eqref{l}), we have, initially for $f\in\gh_{\kappa,\,m}$ but extending to $f\in\gh_\kappa$,
\begin{multline}\label{eq:hardyang}
\left(f,|x|^{-2}f\right)_{L^2\left(\rz^3\,:\,\cz^4\right)}\\
\leq \frac{\left(f^+,\,\left(p_{|\kappa|-\theta(\kappa)}^{(r)}\right)^2f^+\right)_{L^2\left(\rz_+\right)}
}{\left(|\kappa|-\frac12\sgn(\kappa)\right)^2}\,+ \,\frac{\left(f^-,\left(p_{|\kappa|-\theta(-\kappa)}^{(r)}\right)^2f^-\right)_{L^2\left(\rz_+\right)}}{\left(|\kappa|+\frac12\sgn(\kappa)\right)^2}\\
\leq 2\frac{\left(f,p_{\ell_\kappa}^2 f\right)_{L^2\left(\rz^3\,:\,\cz^4\right)}}{
\kappa^2}.
\end{multline}


Since there exist $d\in\rz_+$ such that for all $b\in\rz_+$ and $\ell\in\nz_0$
\[
\left\|\left(p_\ell^{(r)}+b\right)\left(C_\ell^{(r)}+b\right)^{-1}\right\|_{L^2\left(\rz_+,\dr\right)}
\leq d\left(b^{-1/2}\one_{\left\{b\,\leq\,1\right\}} + \one_{\left\{b\,>\,1\right\}}\right)
\]
(Frank et al. \cite[Formula~(5.7)]{Franketal2020P}), this implies with $b:=a_\gamma\kappa^{-2}$ and $\ell:=\ell_\kappa$
\begin{equation}\label{eq:plancherelang}
\left\|\left(p_{\ell_\kappa}+a_\gamma \kappa^{-2}\right)\left(C_{\ell_\kappa}+a_\gamma
\kappa^{-2}\right)^{-1}\right\|_{\gh_{\kappa}} \leq d\left(\frac{|\kappa|}{\sqrt{a_\gamma}}\one_{\left\{a_\gamma\,\leq\,
\kappa^2\right\}} + \one_{\left\{a_\gamma\,>\,\kappa^2\right\}}\right).
\end{equation}
We claim that $K_\gamma:=\lceil \sqrt{a_\gamma}\rceil+1$ and $a_\gamma:=\max\{a_{\gamma,\,C},2d^2\}$ are constants that have the claimed properties: the triangle inequality, Lemma~\ref{boundkineticenergy}, and the estimates~\eqref{eq:hardyang} and~\eqref{eq:plancherelang} imply for $f\in\ghpk$
\begin{equation}\label{eq:dcc}
\begin{split}
\left\|\left(F_\gamma+\frac{a_\gamma}{\kappa^2}\right)f\right\| &=
\left\|\left(D_0-1-\frac\gamma{|x|}+\frac{a_\gamma}{\kappa^2}\right)f\right\|\\
&\geq \left\|\left(D_0-1+\frac{a_\gamma}{\kappa^2}\right)f\right\|-\gamma\left\||x|^{-1}f\right\|\\
&\geq \left(1-\gamma\sqrt{2/a_\gamma}d\right)
\left\|\left(C_{\ell_\kappa}+\frac{a_\gamma}{\kappa^2}\right)f\right\|\\
&\geq (1-\gamma)
\left\|\left(C_{\ell_\kappa}+\frac{a_\gamma}{\kappa^2}\right) f\right\|,
\end{split}
\end{equation}
by definition of $a_\gamma$ and $K_\gamma$. Thus,
\[
\left(C_{\ell_\kappa}\left|_\ghpk+a_\gamma \kappa^{-2}\right)\right)^2
\leq (1-\gamma)^{-2}\left(f_{\gamma,\,\kappa}+a_\gamma \kappa^{-2}\right)^2.
\]
By operator monotonicity of the square root and since $f_{\gamma,\,\kappa}+a_\gamma \kappa^{-2}\geq0$ the last bound implies
\begin{equation}\label{eq:boundChandraDirac}
\left.C_{\ell_\kappa}\right|_\ghpk+\frac{a_\gamma}{\kappa^2}\leq (1-\gamma)^{-1}\left(f_{\gamma,\,\kappa}+\frac{a_\gamma}{\kappa^2}\right)\leq (1-\gamma)^{-1}\left(f_{0,\,\kappa}(V)+\frac{a_\gamma}{\kappa^2}\right).
\end{equation}
Next, \eqref{eq:hardyang} and~\eqref{eq:plancherelang} allow us to estimate
\begin{equation}
U_1 \leq \|rU_1\|_\infty \left(C_{\ell_\kappa}+a_\gamma \kappa^{-2}\right).
\end{equation}
Moreover, by~\eqref{eq:genSobolevDual} and the definition of $\cK_{s,\,0}$
\[
|U_2|^{2s}
\leq A_{s,\,a_\gamma}\left\||U_2|^{2s}\right\|_{\cK_{s,\,0}}\left(C_{\ell_\kappa} +a_\gamma \kappa^{-2}\right)^{2s}
\]
Thus, by operator monotonicity of $x\mapsto x^s$ with $s\in(0,1]$,
\[
U \leq
\left[\left\|rU_1\right\|_\infty+A_{s,\,a_\gamma}^{1/(2s)}\left\||U_2|^{2s}\right\|_{\cK_{s,\,0}}^{1/(2s)}\right] \left(C_{\ell_\kappa}+a_\gamma \kappa^{-2}\right).
\]
Combining this bound with~\eqref{eq:boundChandraDirac}, we obtain for sufficiently small $|\lambda|$,
\[
f_{0,\,\kappa}(V)+a_\gamma \kappa^{-2}
\geq (1-\gamma)\left(C_{\ell_\kappa}+a_\gamma \kappa^{-2}\right)
\geq \lambda U,
\]
thereby proving the assertion.
\end{proof}

We are now ready to prove Proposition~\ref{dominantsigma}.

\begin{proof}[Proof of Proposition~\ref{dominantsigma}]
Let $d_{\kappa,\,\lambda}$ denote the orthogonal projection onto the negative spectral subspace of $F_0(V+\lambda U)$ in $\ghpk$. By the variational principle, we obtain
\begin{equation}
s_{\kappa,\,\lambda} := \tr_\kappa F_0(V+\lambda U)_- - \tr_\kappa F_0(V)_-
\leq \lambda\tr\left(d_{\kappa,\lambda}U\right)
\end{equation}
Similar to~\cite[Equation~(19)]{Iantchenkoetal1996} we set
\begin{equation}
\begin{split}
A :=\,&d_{\kappa,\,\lambda}\left(f_{0,\,\kappa}\left(V+\lambda U\right)+b_\kappa\right)^s,\\
B :=\,&\left(f_{0,\,\kappa}\left(V+\lambda U\right)+b_\kappa\right)^{-s}\Lambda_\gamma\left(C_{\ell_\kappa}+b_\kappa\right)^s,\\
C :=\,&\left(C_{\ell_\kappa}+b_\kappa\right)^{-s} U \left(C_{\ell_\kappa}+b_\kappa\right)^{-s}
\end{split}
\end{equation}
yielding
\[
s_{\kappa,\,\lambda}\leq\lambda \tr \left(ABCB^*A^*\right).
\]
We choose $ b_\kappa:=b_\gamma/\kappa^2$ with some sufficiently large $b_\gamma$, that is going to be determined later. We start by estimating $\|A\|$ using Lemma~\ref{boundA} which is applicable since $U_2^{2s}\in\cK_{s',\,4(s-s')}\subseteq\cK_{s,\,0}$ by Lemma~\ref{inclusions}. Since $d_{\kappa,\,\lambda}$ projects onto the negative spectral subspace of $F_0(V+\lambda U)$ on $\ghpk$, Lemma~\ref{boundA} implies that there are $\lambda_2>0$ and $K_\gamma\in\nz$ such that for all $\lambda\in\rz$ with $|\lambda|<\lambda_2$ and all $\kappa\in\zp$ with $|\kappa|\geq K_\gamma$, we have $f_{0,\kappa}(V+\lambda U)+b_\gamma\kappa^{-2}\geq(b_\gamma-a_\gamma)
\kappa^{-2}$ which is strictly positive for $b_\gamma>a_\gamma$ which we will assume from now on. In particular $\|A\|\leq b_\gamma^{s} |\kappa|^{-2s}$.

Next, $\|C\|_{\gS^1(\gh_{\kappa})} \leq A_{s,\,b_\gamma}|\kappa|\|U\|_{\cK_{s,\,0}}$ is an immediate consequence of Lemma~\ref{genReltrclass}.

We now show the boundedness of $B$. We write $B=B_1B_2$ where
\begin{align*}
B_1 & := \left(f_{0,\,\kappa}\left(V+\lambda U\right)+b_\gamma\kappa^{-2}\right)^{-s} \left(f_{0,\,\kappa}\left(V+\lambda U_1\right)+b_\gamma\kappa^{-2}\right)^s\\
B_2 & := \left(f_{0,\,\kappa}\left(V+\lambda U_1\right)+b_\gamma\kappa^{-2}\right)^{-s}\Lambda_\gamma\left(C_{\ell_\kappa}+b_\gamma\kappa^{-2}\right)^s
\end{align*}
as operators in $\ghpk$. To estimate $\|B_2\|$, we wish to show
\begin{equation}\label{eq:boundB2squared}
\ppp\left(C_{\ell_\kappa}+b_\gamma\kappa^{-2}\right)^{2}\ppp
\leq 4\left(f_{0,\,\kappa}\left(V+\lambda U_1\right)+b_\gamma\kappa^{-2}\right)^{2}.
\end{equation}
Believing this estimate for the moment, we can use the operator monotonicity of $x\mapsto x^{s}$ with $s\in(0,1]$ and Frank and Geisinger's inequality~\eqref{eq:FG}. For $A=(C_{\ell_\kappa}+b_\gamma\kappa^{-2})^2$ in $\ghpk$, $P=\Lambda_\gamma$, $f(x)=x^s$, and $0<s\leq1$, formula~\eqref{eq:FG} reads in $\ghpk$
\begin{multline}\label{eq:SDgen}
\Lambda_\gamma \left(C_{\ell_\kappa}+\frac{b_\gamma}{\kappa^2}\right)^{2s}\Lambda_\gamma
\leq \Lambda_\gamma\left(\Lambda_\gamma\left(C_{\ell_\kappa}+\frac{b_\gamma}{\kappa^2}\right)^2\Lambda_\gamma\right)^s\Lambda_\gamma \\= \left(\Lambda_\gamma\left(C_{\ell_\kappa}+\frac{b_\gamma}{\kappa^2}\right)^2\Lambda_\gamma\right)^s.
\end{multline}
Combining this inequality with~\eqref{eq:boundB2squared} would establish the boundedness of $B_2$.

To prove~\eqref{eq:boundB2squared}, we use the inverse inequality and $\Lambda_\gamma\leq1$ and estimate for $f\in\ghpk$
\begin{multline}\label{eq:boundU1quad1}
\left\|\left(F_0\left(V+\lambda U_1\right)+\frac{b_\gamma}{\kappa^2}\right)f\right\|\\
= \left\|\Lambda_\gamma\left(D_0-\gamma/r+\gamma/r-V-\lambda U_1-1+\frac{b_\gamma}{\kappa^2}\right)\Lambda_\gamma f\right\|\\
\geq \left\|\Lambda_\gamma\left(D_\gamma-1+b_\gamma\kappa^{-2}\right)\Lambda_\gamma f\right\| - \left(\gamma+|\lambda|\left\|rU_1\right\|_\infty\right)\left\||x|^{-1}\Lambda_\gamma f\right\|.
\end{multline}
Using that $\Lambda_\gamma$ and $D_\gamma-1+b_\gamma\kappa^{-2}$ commute and Lemma~\ref{boundkineticenergy}, we estimate further
\begin{multline}\label{eq:boundU1quad2}
\left\|\Lambda_\gamma\left(D_\gamma-1+b_\gamma\kappa^{-2}\right)\Lambda_\gamma f\right\| =\left\|\left(D_\gamma-1+b_\gamma\kappa^{-2}\right)\Lambda_\gamma f\right\|\\
\mkern-10mu
\geq \left\|\left(D_0-1+b_\gamma\kappa^{-2}\right)\Lambda_\gamma f\right\| - \gamma\left\||x|^{-1}\Lambda_\gamma f\right\|
\geq \left\|\left(C_{\ell_\kappa}+b_\gamma\kappa^{-2}\right)\Lambda_\gamma f\left\| - \gamma\right\||x|^{-1}\Lambda_\gamma f\right\|.
\end{multline}
Combining~\eqref{eq:boundU1quad1} and~\eqref{eq:boundU1quad2} with~\eqref{eq:hardyang} and~\eqref{eq:plancherelang}, we obtain
\begin{multline*}
\left\|\Lambda_\gamma\left(D_0\left(V+\lambda U_1\right)-1+b_\gamma\kappa^{-2}\right)\Lambda_\gamma f\right\| \\
\mkern-20mu \geq\left[1-\sqrt{2}d\left(2\gamma+\lambda\left\|rU_1\right\|_\infty\right)\left(b_\gamma^{-\frac12}\one_{\left\{b_\gamma\,\leq\,\kappa^2\right\}} + |\kappa|^{-1}\one_{\left\{b_\gamma\,\geq\,\kappa^2\right\}}\right)\right] \left\|\left(C_{\ell_\kappa}+b_\gamma\kappa^{-2}\right)\Lambda_\gamma f\right\|.
\end{multline*}
Choosing
\begin{equation}\label{eq:gammacond}
b_\gamma = 2\max\left\{a_\gamma,8d^2\left(2\gamma+\lambda_2\left\|rU_1\right\|_\infty\right)^2\right\}
\quad\text{and}\quad K_\gamma = \left\lceil\sqrt{b_\gamma}\right\rceil
\end{equation}
with $a_\gamma$ as in Lemma~\ref{boundA} shows
\begin{equation}\label{eq:boundU1quad}
\left(f_{0,\,\kappa}\left(V+\lambda U_1\right)+b_\gamma\kappa^{-2}\right)^2
\geq \frac14 \Lambda_\gamma\left(C_{\ell_\kappa}+b_\gamma\kappa^{-2}\right)^2\big|_\ghpk
\end{equation}
for all $\nz\ni|\kappa|\geq K_\gamma$ and $|\lambda|<\lambda_2$, thereby establishing~\eqref{eq:boundB2squared}. Using~\eqref{eq:boundB2squared}, operator monotonicity of $x\mapsto x^s$ for $s\in(0,1]$, and~\eqref{eq:SDgen}, we eventually obtain
\begin{equation}\label{eq:boundB2}
\Lambda_\gamma\left(C_{\ell_\kappa}+b_\gamma\kappa^{-2}\right)^{2s}\big|_\ghpk
\leq 4^s\left(f_{0,\,\kappa}\left(V+\lambda U_1\right)+b_\gamma\kappa^{-2}\right)^{2s}
\end{equation}
for all $|\kappa|\geq K_\gamma$ and $|\lambda|<\lambda_2$. This shows $\|B_2\|<2^s$.

Now, we turn to $B_1$ and show
\begin{equation}\label{eq:boundB1}
\left(f_{0,\,\kappa}\left(V+\lambda U_1\right)+b_\gamma\kappa^{-2}\right)^{2s}\\
\leq 2\left(f_{0,\,\kappa}\left(V+\lambda U_1\right)+b_\gamma\kappa^{-2}-\lambda \ppp U_2\ppp\right)^{2s}.
\end{equation}
By~\cite[Lemma~3.4]{Franketal2020P}, which we recall in Lemma~\ref{apriori}, Estimate~\eqref{eq:boundB1} holds, provided we can show
\begin{equation}\label{eq:Neidhardt}
\left\|\big|\lambda \Lambda_\gamma U_2\Lambda_\gamma\big|^{s}\left(D_0(V+\lambda U_1)\middle|_{\ghpk}-1+\frac{b_\gamma}{2\kappa^2}\right)^{-s'}\right\|
\leq A_{s,\,s'}\left(\frac{b_\gamma}{2\kappa^2}\right)^{s-s'}
\end{equation}
for a certain constant $A_{s,\,s'}$ and some $1/2<s'<s$. To show~\eqref{eq:Neidhardt}, we first use the Davis--Sherman inequality~\eqref{eq:SD} and~\eqref{eq:genSobolevDual} and obtain
\begin{align*}
\left|\lambda \Lambda_\gamma U_2\Lambda_\gamma\right|^{2s}
&\leq |\lambda|^{2s} \Lambda_\gamma U_2^{2s}\Lambda_\gamma\\
&\leq 4^{s'}A_{s',\,b_\gamma} |\lambda|^{2s} |\kappa|^{-4(s-s')}\left\|U_2^{2s}\right\|_{\cK_{s',\,4\left(s-s'\right)}}\Lambda_\gamma\left(C_{\ell_\kappa}+b_\gamma\kappa^{-2}\right)^{2s'}\Lambda_\gamma
\;\text{ in } \gh_{\kappa}.
\end{align*}
Combining this estimate with~\eqref{eq:boundB2} with $s$ replaced by $s'$, i.e.,
\[
\left\|\left(C_{\ell_\kappa}+b_\gamma\kappa^{-2}\right)^{s'}\Lambda_\gamma\left(D_0\left(V+\lambda U_1\right)\middle|_\ghpk-1+b_\gamma\kappa^{-2}/2\right)^{-s'}\right\| \leq 2^{s'}
\]
shows that the left side of~\eqref{eq:Neidhardt} is bounded by $4^{s'}A_{s',\,b_\gamma}^{1/2}|\lambda|^s\|U_2^{2s}\|_{\cK_{s',\,4(s-s')}}^{1/2} |\kappa|^{-2(s-s')}$.

Thus, there is a $\lambda_3>0$ such that~\eqref{eq:Neidhardt} holds for all $|\lambda|<\lambda_3$ which shows $\|B\|\leq A_{s}$, uniformly in $\lambda$ and $\kappa$.

Combining the bounds on $\|A\|^2$, $\|B\|^2$, and $\|C\|_1$, we find for $|\lambda|<\min\{\lambda_2,\lambda_3\}$ and all $|\kappa|\geq K_\gamma$,
\[
s_{\kappa,\,\lambda}
\leq A_{\gamma,\,s}\lambda \|U\|_{\cK_{s,\,0}} |\kappa|^{1-4s}
\]
what was claimed since $s>1/2$.
\end{proof}



\subsection{Proof of Theorem~\ref{existencerhoh} on the existence of \texorpdfstring{$\rho^H$}{rhoH}}
We will now prove the pointwise bounds on $\rho_{\kappa}^H$ of Theorem~\ref{existencerhoh}. The strategy of the proof is similar to the one of Proposition~\ref{dominantsigma}.

\begin{proof}[{Proof of Theorem~\ref{existencerhoh}}]
Let $d_\kappa$ denote the orthogonal projection onto the negative spectral subspace of $F_\gamma$ in $\ghpk$. Then
\[
\rho_{\kappa}^H(x) = \tr d_\kappa \delta^{(s)}_{|x|} = \tr ABCB^*A^*
\]
where $\delta^{(s)}_R$ is the delta sphere function with radius $R$, i.e., $\delta^{(s)}_R(y):=\delta(|y|-R)/(4\pi R^2)$ and
\begin{equation}
\begin{split}
A:=\,&d_\kappa\left(f_{\gamma,\,\kappa}+\tilde a_\kappa\right)^s,\\
B:=\,&\left(f_{\gamma,\,\kappa}+\tilde a_\kappa\right)^{-s}\Lambda_\gamma\left(C_{\ell_\kappa}+\tilde a_\kappa\right)^s,\\\
C:=\,&\left(C_{\ell_\kappa}+\tilde a_\kappa\right)^{-s}\ \delta_{|x|}^{(s)} \left(C_{\ell_\kappa}+\tilde a_\kappa\right)^{-s}
\end{split}
\end{equation}
with $\tilde a_\kappa:=a_{\gamma,\,C}\kappa^{-2}$, $a_{\gamma,\,C}$ as defined in the beginning of the proof of Lemma~\ref{boundA}. Moreover, the parameter $s$ obeys $1/2<s<3/2-\sigma_\gamma$ and $s\leq3/4$.

First, we have $\|A\|^2\leq a_{s,\,\gamma} |\kappa|^{-4s}$ by~\eqref{eq:eigenvalue}.

Next,
\[
\tr C=\frac{4|\kappa|}{|x|^2}\left[\left(C_{\ell_\kappa}^{(r)}+\tilde a_\kappa\right)^{-2s}(|x|,|x|) + \left(C_{|\kappa|-\theta(-\kappa)}^{(r)}+\tilde a_\kappa\right)^{-2s}(|x|,|x|)\right].
\]
Here, it is crucial to have $s>\tfrac12$, since $\delta_r$ on $\rz_+$ is not form bounded with respect to $(C_{\ell_\kappa}^{(r)})^{2s}$ for any $s\leq\tfrac12$. The diagonal was estimated in~\cite[Lemma~A.2]{Franketal2020P} for $s\in(\tfrac12,\tfrac34]$, namely
\begin{multline*}
\left(C_{\ell_\kappa}^{(r)}+\tilde a_\kappa\right)^{-2s}(r,r) \\
\leq A_{s,\,a_\gamma}\left[\left(\frac{r}{|\kappa|}\right)^{2s-1}\one_{\left\{r\,\leq\,|\kappa|\right\}} + \left(\frac{r}{|\kappa|}\right)^{4s-1}\one_{\left\{|\kappa|\,\leq\,r\,\leq\,\kappa^2\right\}} + |\kappa|^{4s-1}\one_{\left\{r\,\geq\,\kappa^2\right\}}\right].
\end{multline*}
Repeating this computation for
\[
\left(C_{|\kappa|-\theta(-\kappa)}^{(r)}+\tilde a_\kappa\right)^{-2s}(r,r)
\]
shows that the same bound holds also in this case, since $|\ell_\kappa-(|\kappa|-\theta(-\kappa))|=1$.

The uniform boundedness of $B$ in $|\kappa|$ was shown in the proof of Proposition~\ref{dominantsigma} for $\gamma<1$ and all $|\kappa|\geq K_\gamma$ where $K_\gamma$ is given in~\eqref{eq:gammacond}.

For $|\kappa|\leq K_\gamma$, the uniformity of estimates on $\|B\|$ with respect to $|\kappa|$ is not crucial, since only a fixed finite number of angular momentum channels is involved. For these $\kappa$, we write
\begin{equation}
B=\left(B_1\circ B_2\circ B_3\right)^*
\end{equation}
where
\begin{align*}
B_1 :=\, &\left(C_{\ell_\kappa}+\tilde a_\kappa\right)^s \left(p_\ell+\tilde a_\kappa\right)^{-s}, \\
B_2 :=\, &\left(p_\ell+\tilde a_\kappa\right)^{s}\left(\left(f_{\gamma,\,\kappa}+1+\tilde a_\kappa\right)\right)^{-s},\\
B_3 :=\, &\left(f_{\gamma,\,\kappa}+1+\tilde a_\kappa\right)^{s}\left(\left(f_{\gamma,\,\kappa}+\tilde a_\kappa\right)\right)^{-s}.
\end{align*}
Clearly, $\|B_1\|\leq A_{s,\,\gamma}$ in each channel. By Corollary~\ref{hardydomcor}, respectively~\eqref{eq:hardydomcorang}, for fixed $\gamma<1$ and $s<\frac32-\sigma_\gamma$ if $\gamma\geq\frac{\sqrt3}2$, we have $\|B_2\|\leq A_{s,\gamma}$ for all $\kappa\in\zp$. By~\eqref{eq:eigenvalue}
\begin{equation}
\left\|\left(f_{\gamma,\,\kappa}+\tilde a_\kappa+1\right)\left(f_{\gamma,\,\kappa}+\tilde a_\kappa\right)^{-1}\right\| \leq 1 + A \kappa^2 \leq A_{K_\gamma}.
\end{equation}
Thus, by operator monotonicity of $x\mapsto x^s$ with $s\in(0,1]$ and the bounds on $B_1,B_2$, and $B_3$, we obtain $ \|B\|\leq A_{s,\gamma}$. Combining the bounds on $A,B$, and $C$, we obtain
\begin{multline*}
\rho_{\kappa}^H(x) \leq A_{s,\gamma} \frac{|\kappa|^{1-4s}}{|x|^2}\\
\times \left[\left(\frac{|x|}{|\kappa|}\right)^{2s-1}\one_{\{|x| \leq |\kappa|\}}+\left(\frac{|x|}{|\kappa|}\right)^{4s-1}\one_{\{|\kappa| \leq |x| \leq |\kappa|^2\}} + |\kappa|^{4s-1}\one_{\{|x| \geq |\kappa|^2\}}\right].
\end{multline*}
In particular, the right side is summable for $s>1/2$ and one finally obtains
\begin{align*}
\rho^H(x) = \sum_{\kappa\,\in\,\zp}^\infty \rho_\kappa^H(x)
\leq A_{s,\,\gamma} \left(|x|^{2s-3}\one_{\left\{|x|\,\leq\,1\right\}} + |x|^{-3/2}\one_{\left\{|x|\,>\,1\right\}}\right).
\end{align*}
Recalling the assumptions on $s$ concludes the proof of Theorem~\ref{existencerhoh}.
\end{proof}

\section{Proof of the strong Scott conjecture}
We are now in position to prove Theorem~\ref{convfixedl}, i.e., the strong Scott conjecture for fixed angular momentum, or to be more accurate, for fixed spin-orbit coupling $\kappa$.

We start the proof with some general remarks. Since the statement is linear with respect to $U$, we can assume without loss of generality that $U$ is non-negative and belongs either to $r^{-1}L^\infty_0 (\rz_+)$ or to $\cD_\gamma^{(0)}$.

Given a spherically symmetric potential $U$ define $U_c$ by $U_c(x):=c^2U(cx)$. Furthermore, using~\eqref{eq:manyfurry} for $N=Z$ and fixing $\kappa\in\zp$, we introduce the quadratic form
\begin{equation}
\begin{split}
\cE_{c,\,Z,\,\lambda,\,\kappa}:\bigwedge_{n=1}^Z\Lambda_{c,\,Z}\left(\cS\left(\rz^3:\cz^4\right)\right) &\rightarrow\cz,\\
\psi &\mapsto \cE_{c,\,Z}[\psi]-\lambda \sum_{\nu=1}^Z\big(\psi,\left(\Pi_\kappa\circ U_c\circ\Pi_\kappa\right)_\nu\psi\big),
\end{split}
\end{equation}
if this form is defined and bounded from below.

If $U\in\cD_\gamma^{(0)}$ then $U^{2s}\in\cK_{s'}^{(0)}$ and thus
\[
\left(\Pi^+_{\kappa}U\Pi_{\kappa}^+\right)^{2s} \leq a_{s,\,s',\,\gamma}\left(f_{\gamma,\,\kappa}+1\right)^{2s'}
\]
for $1/2<s'<s$ by the proof of Proposition~\ref{diffu2}. Thus, by Proposition~\ref{diffgen} and Lemma~\ref{apriori}, $U$ is infinitesimally form bounded with respect to $f_{\gamma,\,\kappa}$. Hence, $\cE_{c,\,Z,\,\lambda,\,\kappa}$ is defined and bounded from below.

If $U\in r^{-1}L^\infty_0 ([0,\infty))$, the same follows from Kato's inequality and Corollary~\ref{hardydomcor} for all $\lambda$ in an $Z$ independent open neighborhood of zero.

We now rewrite the expectation of the one-particle perturbation $\Pi^+_\kappa U_c\Pi^+_\kappa$ in the state $\Lambda$ in terms of its ground state density (see~\eqref{eq:furryangdensity}) $\rho_{\kappa}(x)$ in channel $\kappa$; we have
\begin{equation}\label{eq:4.2}
\int_{\rz^3} \rho_\kappa(x)U_c(x)\,\dx = \frac{1}{D\lambda}\sum_{d=1}^D\left(\cE_{c,\,Z}[\psi_d]-\cE_{c,\,Z,\,\lambda,\,\kappa}[\psi_d]\right).
\end{equation}
It obviously depends only superficially on $\lambda$. To estimate this from above we pick $\lambda>0$, use the upper bound on $\cE_{c,\,Z}$~\cite{HandrekSiedentop2015} (Scott correction for the Furry operator) and a lower bound on $\cE_{c,\,Z,\,\lambda,\,\kappa}$ by the correlation inequality of Mancas et al.~\cite{Mancasetal2004} (MMS). This reduces the problem to a one-particle Furry operator with screened Coulomb potential given by the Thomas--Fermi density. This one-particle problem can be treated by the methods developed in the previous sections. The corresponding lower bound will be for free by reversing the sign of $\lambda$.

We begin with the lower bound on $\cE_{c,\,Z,\,\lambda,\,\kappa}$ recalling a special case of the MMS correlation inequality: we write $\rho_Z^{\TF}$ for the minimizer of the Thomas--Fermi functional for a neutral atom (Lieb and Simon~\cite[Theorem~II.20]{LiebSimon1977}). Next we define a ball centered at $x$ with radius $R_Z(x)$ defined by
\[
\int_{|x-y|\,\leq\,R_Z(x)}\rho_Z^{\TF}(y)\,\dy=\frac12.
\]
The screening potential to be used is $\chi_Z(x):=\int_{|x-y|\,\geq\,R_Z(x)}\rho_Z^{\TF}(y)|x-y|^{-1}\,\dy$. Then~\cite{Mancasetal2004} asserts
\begin{equation}\label{MMS}
\sum_{1\,\leq\,\nu\,<\,\mu\,\leq\,Z}|x_\nu-x_\mu|^{-1}
\geq \sum_{\nu=1}^Z\chi_Z(x_\nu)-\gD\left[\rho_Z^{\TF}\right]
\end{equation}
with the short-hand notation
\begin{align*}
\gD[\rho] := \frac12 \int_{\rz^3\,\times\,\rz^3}\frac{\overline{\rho(x)}\rho(y)}{|x-y|}\dx\,\dy\,.
\end{align*}
This allows to eliminate all two-particle terms in $\cE_{c,\,Z,\,\lambda,\,\kappa}$.


\begin{lemm}\label{correlation}
For sufficiently small $\lambda$ and all $L\in\nz$
\begin{multline}\label{Korrelation}
D^{-1}\sum_{d=1}^D\cE_{c,\,Z,\,\lambda,\,\kappa}\left[\psi_d\right]
\geq -\sum_{|{\kappa'}|\,<\,L}\tr_{\kappa'}\left(F_{c,\,Z}\left(\lambda\Pi_{\kappa} U_c\Pi_{\kappa}\right)_{-} \right)  \\
- \sum_{Z/2\,\geq\,|{\kappa'}|\,\geq\,L}\tr_{\kappa'} F_{c,\,Z}\left(-\chi_Z+\lambda \Pi_\kappa U_c\Pi_\kappa\right)_- - \gD\left[\rho_Z^{\TF}\right]
\end{multline}
\end{lemm}

\begin{proof}
By MMS, using the one-particle density matrix $\gamma_\Lambda$ of $\Lambda$, partial wave analysis and $\chi_Z\geq0$, we have
\begin{multline}
D^{-1}\sum_{d=1}^D\cE_{c,\,Z,\,\lambda,\,\kappa}\left[\psi_d\right] \\
\geq \sum_{|{\kappa'}|\,<\,L}\tr_{\kappa'}\left[\gamma_\Lambda F_{c,\,Z}\left(\lambda\Pi_\kappa U_c\Pi_\kappa\right)\right] + \sum_{|{\kappa'}|\,\geq\,L}\tr_{\kappa'}\left[\gamma_\Lambda F_{c,\,Z}\left(-\chi_Z+\lambda \Pi_\kappa U_c\Pi_\kappa\right)\right] - \gD\left[\rho_Z^{\TF}\right]
\end{multline}
for any $L\in\nz$. Since the energy is increasing in $|\kappa'|$ (see Lemma~\ref{boundA}) and $\tr\gamma_\Lambda=Z$ we can minimize in the one-particle density matrix under this constraint. The resulting summands will surely vanish, if $|{\kappa'}|>\tfrac{Z}{2}$, i.e., we can cut off the series at $\tfrac{Z}{2}$. Dropping the requirement of $Z$ particles only gives the wanted result.
\end{proof}

Next we recall the asymptotic expansion of the ground state energy, i.e.,
\[
D^{-1}\sum_{d=1}^D\cE_{c,\,Z}\left[\psi_d\right] = E^{\TF}(Z) + \left(\frac12-s(\gamma)\right)Z^2+\cO\left(Z^{47/24}\right)
\quad \text{as}\ Z\to\infty
\]
(Handrek and Siedentop~\cite[Theorem~1]{HandrekSiedentop2015}) with the finite spectral shift
\begin{equation}\label{eq:spectralshift}
s(\gamma) =\gamma^{-2}\sum\nolimits_{{\kappa'}\,\in\,\zp}\left[\tr_{\kappa'}\left(F_\gamma\right)_{-} - \left|{\kappa'}\right|\sum\nolimits_{n\,\in\,\nz}\gamma^2\left(n+\ell_{\kappa'}\right)^{-2}\right].
\end{equation}
In fact, with $L:=Z^{1/9}$ their proof gives the stronger chain of inequalities
\begin{multline}\label{scottoperator}
-\const Z^{47/24} + E^{\TF}(Z)+\left(\frac12-s(\gamma)\right)Z^2\\
\begin{aligned}
&\leq -\sum_{\left|{\kappa'}\right|\,<\,L}\tr_{\kappa'} \left(F_{c,\,Z}\right)_{-} - \sum_{Z/2\geq |{\kappa'}|\geq L}\tr_{\kappa'} F_{c,\,Z}\left(-\chi_Z\right)_{-} -\gD\left[\rho_Z^{\TF}\right]\leq D^{-1}\sum_{d=1}^D\cE_{c,\,Z}\left[\psi_d\right] \\
&\leq E^{\TF}(Z)+\left(\frac12-s(\gamma)\right)Z^2+\const Z^{47/24}
\end{aligned}
\end{multline}
implying

\begin{multline}\label{wichtig}
D^{-1}\sum_{d=1}^D\cE_{c,\,Z}\left[\psi_d\right]\\
= -\sum_{\left|{\kappa'}\right|\,<\,Z^\frac19}\tr_{\kappa'} \left(F_{c,\,Z}\right)_{-} - \sum_{Z/2\,\geq\,\left|{\kappa'}\right|\,\geq\,Z^\frac19}\tr_{\kappa'} F_{c,\,Z}\left(-\chi_Z\right)_{-} -
\gD\left[\rho_Z^{\TF}\right] + O\left(Z^\frac{47}{24}\right).
\end{multline}

The above preparations allow us to conclude the
\begin{proof}[Proof of Theorem~\ref{convfixedl}]
We divide~\eqref{eq:4.2} by $c^2$, use the bounds~\eqref{wichtig} and~\eqref{Korrelation}, and rescale $x\to x/c$. We get for positive $\lambda$
\begin{multline}\label{eq:reduced}
\int_{\rz^3}c^{-3}\rho_\kappa(x/c)U(x)\rd x =\frac1{c^2D\lambda}\sum_{d=1}^D\left(\cE_{c,\,Z}\left[\psi_d\right]-\cE_{c,\,Z,\,\lambda,\,\kappa}\left[\psi_d\right]\right) 
\\
\mkern-15mu\leq\frac1\lambda
\!
\begin{cases}
\tr\!\left[f_{\gamma,\,\kappa}\left(\lambda U\right)_{-}-f_{\gamma,\,\kappa}(0)_{-}\right] &\!\!\!\!|\kappa|<Z^\frac19\\[0.25em]
\tr\!\left[f_{\gamma,\,\kappa}\big(-c^{-2}\chi_Z\left(\cdot/c\right)+\lambda U\big)-f_{\gamma,\,\kappa}\big(-c^{-2}\chi_Z\left(\cdot/c\right)\big)_{-}\right]& \!\!\!\!|\kappa|\geq Z^\frac19
\end{cases} + \const Z^{-\frac1{24}}
\end{multline}
where we use that $c=\gamma^{-1}Z$. Taking $c$ to $\infty$ gives
\begin{equation}\label{oben}
\limsup_{c\,\to\,\infty}\int_{\rz^3}c^{-3}\rho_\kappa(x/c)U(x)\rd x
\leq \frac{\tr f_{\gamma,\,\kappa}(\lambda U)_--\tr f_{\gamma,\,\kappa}(0)_-}{\lambda}.
\end{equation}
Taking $\lambda<0$ gives the reverse inequality
\begin{equation}\label{unten}
\liminf_{c\,\to\,\infty}\int_{\rz^3}c^{-3}\rho_\kappa(x/c)U(x)\rd x
\geq \frac{\tr f_{\gamma,\,\kappa}(\lambda U)_--\tr f_{\gamma,\,\kappa}(0)_-}{\lambda}.
\end{equation}
Propositions~\ref{diffu2} and~\ref{diffu1} show that the right sides of~\eqref{oben} and~\eqref{unten} tend to $\int_{\rz^3} \rho_{\kappa}^H(x)U(|x|)\,\dx$ thus yielding the existence of the limit and its limit quod erat demonstrandum.
\end{proof}


The proof of Theorem~\ref{convalll} is standard and completely analogous to that of Frank et al. \cite[Theorem~1.2]{Franketal2020P} (see also Iantchenko et al. \cite[Theorem~2]{Iantchenkoetal1996}). We therefore merely sketch the argument and refer to~\cite[Section~1.3]{Franketal2020P} for a detailed exposition: The assertion follows from the proof of Theorem~\ref{convfixedl} by summing over $\kappa$ provided we can interchange the summation and the limits $Z\to\infty$ and $\lambda\to0$. This, however is secured by Proposition~\ref{dominantsigma} which allows to apply the Weierstra\ss~criterion for uniform converge.




\appendix

\section{Partial Wave Analysis}
\label{Anhang1}
\label{a:partialwaves}
We collect some notations and known facts about the partial wave analysis of Dirac operators (see, e.g., Evans et al. \cite{Evansetal1996}, Balinsky and Evans~\cite[Section~2.1]{BalinskyEvans2011}, and Thaller~\cite[Sections~4.6.3-4.6.5]{Thaller1992}).

Let $Y_{\ell,\,m}$ be the spherical harmonics on the unit sphere $\bbS^2$ obeying the normalization condition $\int_{\bbS^2}|Y_{\ell,\,m}|^2\,\domega=1$ where $\domega$ is the usual surface measure on $\bbS^2$. If $|m|>\ell$, we set $Y_{\ell,\,m}\equiv0$. We begin by observing that those of the spherical spinors

\begin{equation}\label{2.6}
\Omega_{\ell,\,m,\,s}(\omega) :=
\begin{pmatrix} 2s\sqrt{\frac{\ell+\frac12+2sm}{2\ell+1}}
Y_{\ell,\,m-\frac12}(\omega)\\ \sqrt{\frac{\ell+\frac12-2sm}{2\ell+1}} Y_{\ell,\,m+\frac12}(\omega)
\end{pmatrix}
\end{equation}
with $\ell=0,1,2,\,\dots$ and $m=-\ell-\frac12,\,\dots,\,\ell+\frac12$, that do not vanish, form an orthonormal basis of $L^2(\bbS^2:\cz^2)$ (see, e.g., Evans et al. \cite[Equation~(7)]{Evansetal1996}).

Moreover, they are joint eigenfunctions of $L^2$, $J^2$ ($J=L+S$ being the total angular momentum), and $J_3$ with respective eigenvalues $\ell(\ell+1)$, $(\ell+s)(\ell+s+1)$, and $m$.

Introducing the spin-orbit operator $K=\beta(J^2-L^2+1/4)$, there is an orthonormal basis of eigenvectors $\Phi^\sigma_{\kappa,\,m}$ of $L^2(\bbS^2:\cz^4)$ such that $J^2\Phi^\sigma_{\kappa,\,m}=j_\kappa(j_\kappa+1)\Phi^\sigma_{\kappa,\,m}$, $J_3\Phi^\sigma_{\kappa,\,m}=m\Phi^\sigma_{\kappa,\,m}$, and $K\Phi^\sigma_{\kappa,\,m}=\kappa\Phi^\sigma_{\kappa,\,m}$ with $j_\kappa:=|\kappa|-1/2$ introduced in~\eqref{l}, $m\in\{-j_\kappa,\,\dots,\,j_\kappa\}$, $\kappa\in\zp$, and $\sigma\in\{+,-\}$. A standard choice is
\begin{equation}\label{fi}
\Phi^+_{\kappa,\,m} := \left(
\begin{matrix}
\ri\sgn(\kappa)\Omega_{\ell_\kappa,\,m,\,\frac12\sgn(\kappa)}\\
0
\end{matrix}
\right), \
\Phi^-_{\kappa,\,m}:=\left(
\begin{matrix}
0\\
-\sgn(\kappa)\Omega_{\ell_\kappa+\sgn(\kappa),\,m,\,-\frac12\sgn(\kappa)}
\end{matrix}
\right).
\end{equation}
Using these spinors, we introduce
\begin{multline}\label{a3}
\gh_{\kappa,\,m}\\
 :=\spanl\left\{x\mapsto \frac{f^+(|x|)}{|x|} \Phi_{\kappa,\,m}^+\left(\frac x{|x|}\right) + \frac{f^-(|x|)}{|x|}\Phi_{\kappa,\,m}^-\left(\frac x{|x|}\right):f^+,f^-\in L^2\left(\rz_+\right)\right\},
\end{multline}
\begin{equation}\label{a2}
\gh_\kappa:= \bigoplus_{m=-j_\kappa}^{j_\kappa}\gh_{\kappa,m},\ \gh_\kappa^+:=\Lambda_\gamma\gh_\kappa
\end{equation}
These spaces form an orthogonal decomposition of $L^2(\rz^3:\cz^4)$ and $\Lambda_\gamma(L^2(\rz^3:\cz^4))$.

We write $\Pi_\kappa$, $\Pi_{\kappa,\,m}$, and $\Pi_\kappa^+$ for the orthogonal projection onto $\gh_{\kappa}$, $\gh_{\kappa,\,m}$, and $\ghpk$. If we write -- in abuse of notation -- $\Phi_{\kappa,\,m}^\pm(\omega,\tau)$ for the $\tau^{\rm th}$ component of $\Phi_{\kappa,\,m}^\pm(\omega)$, we can write the action of $\Pi_\kappa$ on $g\in L^2(\rz^3\otimes\{1,\,\dots,\,4\})$ more explicitly as
\begin{equation}\label{a1}
\left(\Pi_\kappa g\right)\left(r\omega,\tau\right) = \sum_{\sigma\,\in\,\{+,\,-\}}\;\sum_{m=-j_\kappa}^{j_\kappa} \Phi^\sigma_{\kappa,\,m}\left(\omega,\tau\right)
\sum_{\tau'=1}^4\int_{\bbS^2}\rd \omega'\, \overline{ \Phi^\sigma_{\kappa,\,m}\left(\omega',\tau'\right)}g\left(r\omega',\tau'\right)
\end{equation}
writing $x=r\omega$ with $r:=|x|$ and $\omega:=x/r$.

Note that Dirac operators with spherical potentials leave the space $\gh_{\kappa,\,m}$ invariant which can be seen explicitly in~\eqref{eq:radialdirac}. Moreover, their eigenvalues depend on $\kappa$ only.

Furthermore note, that $\Pi^+_\kappa$ is also an orthogonal projection, since $\Pi_\kappa$ commutes with $\Lambda_\gamma$ (see~\cite[Equation~(27)]{HandrekSiedentop2015}).

\section{Test function spaces}\label{s:testfunctions}


The test functions for which we prove the strong Scott conjecture belong to the function spaces $\cK_{s}^{(0)}$ and $\cK_{s,\,\delta}$ which were already introduced in Frank et al. \cite{Franketal2020P} and are defined as

\begin{equation}\label{eq:defsigmadeltanod}
\begin{split}
\cK_{s}^{(0)} & :=\left\{W\in L^1_{\loc}\left(\rz_+\right):\|W\|_{\cK_{s}^{(0)}}<\infty\right\}\\
\|W\|_{\cK_{s}^{(0)}} & := \int_0^1 r^{2s-1}|W(r)|\,\dr + \int_1^\infty |W(r)|\,\dr
\end{split}
\end{equation}
and
\begin{equation}\label{eq:defsigmadelta}
\begin{split}
\cK_{s,\,\delta} :=& \left\{W\in L^1_{\loc}\left(\rz_+\right):\|W\|_{\cK_{s,\,\delta}}<\infty\right\}\\
\|W\|_{\cK_{s,\,\delta}} :=& \sup_{R\,\geq\,1} R^\delta\left[\int_0^R \left(\frac{r}{R}\right)^{2s-1}|W(r)|\,\dr + \int_R^{R^2}\left(\frac{r}{R}\right)^{4s-1} |W(r)|\,\dr\right.\\
&\mkern236mu\left.+ R^{4s-1}\int_{R^2}^\infty |W(r)|\,\dr\right]
\end{split}
\end{equation}
for $s\geq1/2$ and $\delta\in[0,2s-1]$. Here, $L^p_{\loc}(\rz_+)$ denotes the space of all functions that belong to $L^p$ on any compact subset of $\rz_+$. We note some basic inclusion properties which already occurred implicitly in~\cite{Franketal2020P}.
\begin{lemm}\label{inclusions}
Let $1/2\leq s'<s$ and $\delta\in[0,2s-1]$. Then the spaces $\cK_{s}^{(0)}$ and $\cK_{s,\delta}$ obey the following inclusion properties.
\begin{enumerate}
\item One has $\cK_{s'}^{(0)}\subseteq\cK_{s}^{(0)}$.

\item One has $\cK_{s,\,\delta}\subseteq\cK_{s,\,0}\subseteq\cK_{s}^{(0)}$.

\item One has $\cK_{s',\,4(s-s')} \subseteq \cK_{s,\,0}$, if additionally $1/2<2s/3+1/6 \leq s'<s$.
\end{enumerate}
\end{lemm}

This means that functions must be smoother at the origin the smaller $s$ is. Moreover, functions belonging to $\cK_{s,\,\delta}$ must decay faster at infinity than those belonging to $\cK_{s}^{(0)}$.

To give a digestible representation of our convergence results, we introduce the test function spaces
\begin{align}\label{eq:deftestnod0}
\begin{split}
\cD_\gamma^{(0)} =
\begin{cases}
& \left\{ W\in\cK_{s}^{(0)}:\; |W|^{2s}\in\cK_{s'}^{(0)} \; \text{for some}\; 1/2<s'<s\leq 1 \right\} \\
& \qquad\qquad\qquad \text{if}\ 0<\gamma<\sqrt3/2.
\\
& \left\{ W\in\cK_{s}^{(0)}:\; |W|^{2s}\in\cK_{s'}^{(0)} \; \text{for some}\; 1/2<s'<s<3/2-\sigma_\gamma \right\} \\
& \qquad\qquad\qquad \text{if}\; \sqrt3/2\leq\gamma<1.
\end{cases}
\end{split}
\end{align}
and
\begin{multline}\label{eq:deftest0}
\cD \\
=\left\{W \in\cK_{s,\,0}\,:\ |W|^{2s}\in\cK_{s',\,4\left(s-s'\right)}\: \text{for some}\: 1/2<2s/3+1/6\leq s'<s\leq 3/4 \right\}.
\end{multline}
We refer to~\cite{Franketal2020P} for an alternative and more convenient representation of the space $\cD_\gamma^{(0)}$ as well as the norm $\|W\|_{\cK_{s,\,\delta}}$ (see their Formulae (34) and (43)). For instance, $r^{-1}L^\infty_0$ functions belong both to $\cK_{s}^{(0)}$ and $\cK_{s,\,\delta}$ for $s>1/2$ and $\delta\in[0,\,2s-1]$. Moreover, one easily verifies $L^1\left(\rz_+\right)\subseteq \cK_{s}^{(0)}$ and $L^1(\rz_+,r^\delta\,\dr)\cap L^1(\rz_+,r^{4s-1+\delta}\,\dr)\subseteq\cK_{s,\,\delta}$ for $s\geq1/2$ and $\delta\in[0,2s-1]$.



\section{Auxiliary tools}

The following lemma, which we quote from~\cite[Lemma~3.4]{Franketal2020P} was inspired by Neidhardt and Zagrebnov~\cite[Lemma~2.2]{NeidhardtZagrebnov1999}.

\begin{lemm}\label{apriori}
Let $A$ be a self-adjoint operator with $\inf\sigma(A)>0$ and let $B$ be an operator which satisfies $B\geq 0$ or $B\leq 0$. Assume that for some numbers $\max\{s',1/2\}<s<1$ one has
\begin{align*}
\left\||B|^s A^{-s'}\right\|<\infty.
\end{align*}
Then $B$ is form bounded with respect to $A$ with relative bound zero. Moreover, if $\||B|^s A^{-s'}\|\leq a_{s,\,s'}M^{s-s'}$ for some constant $a_{s,\,s'}$ depending only on $s$ and $s'$,
\[
\frac12 (A+M)^{2s} \leq (A+B+M)^{2s} \leq 2(A+M)^{2s}.
\]
\end{lemm}



\section*{Acknowledgments}

This research was partly carried out at the Institute for Mathematical Sciences at the National University of Singapore during the program
\emph{Density Functionals for Many-Particle Systems: Mathematical Theory and Physical Applications of Effective Equations}. We are grateful to the IMS and the Julian Schwinger foundation for their hospitality and financial support. Special thanks go to Berthold-Georg Englert who was the heart of the program.


Partial financial support by the Deutsche Forschungsgemeinschaft (DFG, German Research Foundation) through grant SI 348/15-1 (H.S.) and through Germany's Excellence Strategy - EXC-2111 - 390814868 (H.S.) is gratefully acknowledged.


%\newcommand*\cprime{$'$}


\begin{thebibliography}{99}
\bibitem{Baumgartner1976} [Bau76] Baumgartner, Bernhard The Thomas–Fermi-theory as result of a strong-coupling-limit, Commun. Math. Phys., Volume 47 (1976) no. 3, pp. 215-219
\bibitem{BalinskyEvans2011} [BE11] Balinsky, Alexander A.; Evans, William D. Spectral Analysis of Relativistic Operators, World Scientific, 2011
\bibitem{Bethe1933} [Bet33] Bethe, Hans A. Chapter 3, Quantenmechanik der Ein-und Zwei-Elektronenatome, Springer, 1933, pp. 273-560
\bibitem{BurkeGrant1967} [BG67] Burke, V. M.; Grant, Ian P. The effect of relativity on atomic wave functions, Proc. Phys. Soc., Volume 90 (1967) no. 2, pp. 297-314
\bibitem{Bhatia1997} [Bha97] Bhatia, Rajendra Matrix Analysis, Graduate Texts in Mathematics, 169, Springer, 1997
\bibitem{BrownRavenhall1951} [BR51] Brown, Gordon E.; Ravenhall, D. G. On the interaction of two electrons, Proc. R. Soc. Lond., Volume 208 (1951), pp. 552-559
\bibitem{Carlen2010} [Car10] Carlen, Eric Trace inequalities and quantum entropy: an introductory course, Entropy and the Quantum. Arizona School of Analysis With Applications, March 16–20, 2009, University of Arizona, Tucson, AZ, US. (Sims, Robert et al., eds.) (Contemporary Mathematics), Volume 529, American Mathematical Society, 2010, pp. 73-140
\bibitem{Chandrasekhar1931} [Cha31] Chandrasekhar, Subramanyan The maximum mass of ideal white dwarfs, Astrophys. J., Volume 74 (1931), pp. 81-82
\bibitem{Darwin1928} [Dar28] Darwin, Charles G. The wave equation of the electron, Proc. R. Soc. Lond., Volume 118 (1928), pp. 654-680
\bibitem{Davis1957} [Dav57] Davis, Chandler A Schwarz inequality for convex operator functions, Proc. Am. Math. Soc., Volume 8 (1957), pp. 42-44
\bibitem{Dirac1930A} [Dir30] Dirac, Paul A. M. A theory of electrons and protons, Proc. R. Soc. Lond., Volume 126 (1930), pp. 360-365
\bibitem{Evansetal1996} [EPS96] Evans, William D.; Perry, Peter; Siedentop, Heinz The spectrum of relativistic one-electron atoms according to Bethe and Salpeter, Commun. Math. Phys., Volume 178 (1996) no. 3, pp. 733-746
\bibitem{Fermi1927} [Fer27] Fermi, Enrico Un metodo statistico per la determinazione di alcune proprietá dell’atomo, Rend. Accad. Naz. Lincei, Volume 6 (1927) no. 12, pp. 602-607
\bibitem{Fermi1928} [Fer28] Fermi, Enrico Eine statistische Methode zur Bestimmung einiger Eigenschaften des Atoms und ihre Anwendung auf die Theorie des periodischen Systems der Elemente, Z. Phys., Volume 48 (1928), pp. 73-79
\bibitem{FrankGeisinger2016} [FG16] Frank, Rupert L.; Geisinger, Leander Refined semiclassical asymptotics for fractional powers of the Laplace operator, J. Reine Angew. Math., Volume 712 (2016), pp. 1-37
\bibitem{Fournaisetal2020} [FLT20] Fournais, Søren; Lewin, Mathieu; Triay, Arnaud The Scott correction in Dirac–Fock theory, Commun. Math. Phys., Volume 378 (2020) no. 1, pp. 569-600
\bibitem{Franketal2020R} [FMS20] Frank, Rupert L.; Merz, Konstantin; Siedentop, Heinz Relativistic strong Scott conjecture: A short proof (2020) (https://arxiv.org/abs/2009.02474, to appear in the proceedings of “Density Functionals for Many-Particle Systems: Mathematical Theory and Physical Applications of Effective Equations”, September 2-27, 2019, Institute for Mathematical Sciences of the National University of Singapore)
\bibitem{Franketal2021} [FMS21] Frank, Rupert L.; Merz, Konstantin; Siedentop, Heinz Equivalence of Sobolev norms involving generalized Hardy operators, Int. Math. Res. Not., Volume 2021 (2021) no. 3, pp. 2284-2303
\bibitem{Franketal2020P} [FMSS20] Frank, Rupert L.; Merz, Konstantin; Siedentop, Heinz; Simon, Barry Proof of the strong Scott conjecture for Chandrasekhar atoms, Pure Appl. Funct. Anal., Volume 5 (2020) no. 6, pp. 1319-1356
\bibitem{FurryOppenheimer1934} [FO34] Furry, Wendell H.; Oppenheimer, Julius R. On the theory of the electron and positive, Phys. Rev., Volume 45 (1934), pp. 245-262
\bibitem{Franketal2008} [FSW08] Frank, Rupert L.; Siedentop, Heinz; Warzel, Simone The ground state energy of heavy atoms: Relativistic lowering of the leading energy correction, Commun. Math. Phys., Volume 278 (2008) no. 2, pp. 549-566
\bibitem{Franketal2009} [FSW09] Frank, Rupert L.; Siedentop, Heinz; Warzel, Simone The energy of heavy atoms according to Brown and Ravenhall: the Scott correction, Doc. Math., Volume 14 (2009), pp. 463-516
\bibitem{Gordon1928} [Gor28] Gordon, Walter Die Energieniveaus des Wasserstoffatoms nach der Diracschen Quantentheorie, Z. Phys., Volume 48 (1928), pp. 11-14
\bibitem{Heisenberg1925} [Hei25] Heisenberg, Werner Über quantentheoretische Umdeutung kinematischer und mechanischer Beziehungen, Z. Phys., Volume 33 (1925) no. 1, pp. 879-893
\bibitem{HeilmannLieb1995} [HL95] Heilmann, Ole J.; Lieb, Elliott H. The electron density near the nucleus of a large atom, Phys. Rev., Volume 52 (1995) no. 5, pp. 3628-3643
\bibitem{HandrekSiedentop2015} [HS15] Handrek, Michael; Siedentop, Heinz The ground state energy of heavy atoms: the leading correction, Commun. Math. Phys., Volume 339 (2015) no. 2, pp. 589-617
\bibitem{Hughes1986} [Hug86] Hughes, Webster An Atomic Energy Lower Bound that Gives Scott’s Correction, Ph. D. Thesis, University of Princeton, Department of Mathematics, USA (1986)
\bibitem{Hughes1990} [Hug90] Hughes, Webster An atomic lower bound that agrees with Scott’s correction, Adv. Math., Volume 79 (1990) no. 2, pp. 213-270
\bibitem{Iantchenko1997} [Ian97] Iantchenko, Alexei The electron density in intermediate scales, Commun. Math. Phys., Volume 184 (1997) no. 2, pp. 367-385
\bibitem{Iantchenkoetal1996} [ILS96] Iantchenko, Alexei; Lieb, Elliott H.; Siedentop, Heinz Proof of a conjecture about atomic and molecular cores related to Scott’s correction, J. Reine Angew. Math., Volume 472 (1996), pp. 177-195
\bibitem{IantchenkoSiedentop2001} [IS01] Iantchenko, Alexei; Siedentop, Heinz Asymptotic behavior of the one-particle density matrix of atoms at distances $${Z}^ {-1}$$ from the nucleus, Math. Z., Volume 236 (2001) no. 4, pp. 787-796
\bibitem{Ivrii2019} [Ivr19] Ivrii, Victor Strong Scott conjecture (2019) (https://arxiv.org/abs/1908.05478)
\bibitem{Kalf1976} [Kal76] Kalf, Hubert The virial theorem in relativistic quantum mechanics, J. Funct. Anal., Volume 21 (1976) no. 4, pp. 389-396
\bibitem{KlausWust1978} [KW78] Klaus, Martin; Wüst, Rainer Characterization and uniqueness of distinguished self-adjoint extensions of Dirac operators, Commun. Math. Phys., Volume 64 (1978), pp. 171-176
\bibitem{Lieb1980} [Lie79] Lieb, Elliott H., Mathematical Problems in Theoretical Physics. Proceedings of the International Conference on Mathematical Physics. Lausanne 1979 (Lecture Notes in Physics), Volume 116 (1979), pp. 553-569
\bibitem{Lieb1981} [Lie81] Lieb, Elliott H. Thomas–Fermi and related theories of atoms and molecules, Rev. Mod. Phys., Volume 53 (1981) no. 4, pp. 603-641
\bibitem{Lieb1982A} [Lie82] Lieb, Elliott H. Analysis of the Thomas–Fermi–von Weizsäcker equation for an infinite atom without electron repulsion, Commun. Math. Phys., Volume 85 (1982) no. 1, pp. 15-25
\bibitem{LiebLiberman1982} [LL82] Lieb, Elliott H.; Liberman, David A. Numerical calculation of the Thomas–Fermi–von Weizsäcker function for an infinite atom without electron repulsion, 1982 no. LA-9186-MS
\bibitem{LiebSimon1977} [LS77] Lieb, Elliott H.; Simon, Barry The Thomas–Fermi theory of atoms, molecules and solids, Adv. Math., Volume 23 (1977) no. 1, pp. 22-116
\bibitem{Luders1951} [Lüd51] Lüders, Gerhart Über die Zustandsänderung durch den Meßprozeß, Ann. der Physik, VI. F., Volume 8 (1951), pp. 322-328
\bibitem{Mittleman1981} [Mit81] Mittleman, Marvin H. Theory of relativistic effects on atoms: Configuration-space Hamiltonian, Phys. Rev., Volume 24 (1981) no. 3, pp. 1167-1175
\bibitem{MorozovMuller2017} [MM17] Morozov, Sergey; Müller, David Lower bounds on the moduli of three-dimensional Coulomb–Dirac operators via fractional Laplacians with applications, J. Math. Phys., Volume 58 (2017) no. 7, 072302, 22 pages
\bibitem{Mancasetal2004} [MMS04] Mancas, Paul; Klaus Müller, A. M.; Siedentop, Heinz The optimal size of the exchange hole and reduction to one-particle Hamiltonians, Theoretical Chemistry Accounts, Volume 111 (2004) no. 1, pp. 49-53
\bibitem{MatteStockmeyer2010} [MS10] Matte, Oliver; Stockmeyer, Edgardo Spectral theory of no-pair Hamiltonian, Rev. Math. Phys., Volume 22 (2010) no. 1, pp. 1-53
\bibitem{Nenciu1976} [Nen76] Nenciu, Gheorghe Self-adjointness and invariance of the essential spectrum for Dirac operators defined as quadratic forms, Commun. Math. Phys., Volume 48 (1976) no. 3, pp. 235-247
\bibitem{NeidhardtZagrebnov1999} [NZ99] Neidhardt, Hagen; Zagrebnov, Valentin A. Fractional powers of self-adjoint operators and Trotter-Kato product formula, Integral Equations Oper. Theory, Volume 35 (1999) no. 2, pp. 209-231
\bibitem{Oelker2019} [Oel19] Oelker, Martin Johannes On Domain, Self-Adjointness, and Spectrum of Dirac Operators for Two Interacting Particles, Ph. D. Thesis, Fakultät für Mathematik, Informatik und Statistik, Ludwig-Maximilians-Universität München, Deutschland (2019)
\bibitem{Pauli1926} [Pau26] Pauli, Wolfgang Über das Wasserstoffspektrum vom Standpunkt der neuen Quantenmechanik, Z. Phys., Volume 36 (1926) no. 5, pp. 336-363
\bibitem{Pilkuhn2005} [Pil05] Pilkuhn, Hartmut Relativistic Quantum Mechanics, Texts and Monographs in Physics, Springer, 2005
\bibitem{Schmincke1972} [Sch72] Schmincke, Upke-Walther Distinguished selfadjoint extensions of Dirac operators, Math. Z., Volume 129 (1972), pp. 335-349
\bibitem{Schwinger1980} [Sch80] Schwinger, Julian Thomas–Fermi model: The leading correction, Phys. Rev. A, Volume 22 (1980) no. 5, pp. 1827-1832
\bibitem{Scott1952} [Sco52] Scott, J. M. C. The binding energy of the Thomas–Fermi atom, Philos. Mag., Volume 48 (1952) no. 343, pp. 859-867
\bibitem{Simon1984} [Sim84] Simon, Barry Fifteen problems in mathematical physics, Perspectives in Mathematics, Birkhäuser, 1984
\bibitem{Sommerfeld1916} [Som16] Sommerfeld, Arnold Zur Quantentheorie der Spektrallinien, Ann. Phys. (Berlin), Volume 356 (1916) no. 17, pp. 1-94
\bibitem{Solovejetal2008} [SSS10] Solovej, Jan Philip; Sørensen, Thomas Østergaard; Spitzer, Wolfgang L. Relativistic Scott correction for atoms and molecules, Commun. Pure Appl. Math., Volume 63 (2010) no. 1, pp. 39-118
\bibitem{Sucher1980} [Suc80] Sucher, Joseph Foundations of the relativistic theory of many-electron atoms, Phys. Rev. A, Volume 22 (1980) no. 2, pp. 348-362
\bibitem{Sucher1984} [Suc84] Sucher, Joseph Foundations of the relativistic theory of many-electron bound states, International Journal of Quantum Chemistry, Volume 25 (1984) no. 1, pp. 3-21
\bibitem{Sucher1987} [Suc87] Sucher, Joseph Relativistic many-electron Hamiltonians, Phys. Scr., Volume 36 (1987) no. 2, pp. 271-281
\bibitem{SiedentopWeikard1986} [SW86] Siedentop, Heinz K. H.; Weikard, Rudi On the leading energy correction for the statistical model of the atom: Non-interacting case, Abh. Braunschw. Wiss. Ges., Volume 38 (1986), pp. 145-158
\bibitem{SiedentopWeikard1987O} [SW87a] Siedentop, Heinz; Weikard, Rudi On the leading energy correction for the statistical model of the atom: Interacting case, Commun. Math. Phys., Volume 112 (1987) no. 3, pp. 471-490
\bibitem{SiedentopWeikard1987U} [SW87b] Siedentop, Heinz; Weikard, Rudi Upper bound on the ground state energy of atoms that proves Scott’s conjecture, Phys. Lett., Volume 120 (1987) no. 7, pp. 341-342
\bibitem{SiedentopWeikard1988} [SW88] Siedentop, Heinz; Weikard, Rudi On the leading energy correction of the statistical atom: Lower bound, Eur. Phys. Lett., Volume 6 (1988) no. 3, pp. 189-192
\bibitem{SiedentopWeikard1989} [SW89] Siedentop, Heinz; Weikard, Rudi On the leading correction of the Thomas–Fermi model: Lower bound, Invent. Math., Volume 97 (1989) no. 1, pp. 159-193 (with an appendix by A. M. K. Müller)
\bibitem{Thaller1992} [Tha92] Thaller, Bernd The Dirac Equation, Texts and Monographs in Physics, Springer, 1992
\bibitem{Thomas1927} [Tho27] Thomas, Llewellyn H. The calculation of atomic fields, Proc. Camb. Philos. Soc., Volume 23 (1927), pp. 542-548
\bibitem{Weizsacker1935} [Wei35] Weizsäcker, C. F. v. Zur Theorie der Kernmassen, Z. Phys., Volume 96 (1935), pp. 431-458
\bibitem{Wust1975} [Wüs75] Wüst, Rainer Distinguished self-adjoint extensions of Dirac operators constructed by means of cut-off potentials, Math. Z., Volume 141 (1975), pp. 93-98
\bibitem{YoneiTomishima1965} [YT65] Yonei, Katsumi; Tomishima, Yasuo On the Weizsäcker correction to the Thomas–Fermi theory of the atom, J. Phys. Soc. Japan, Volume 20 (1965) no. 6, pp. 1051-1057
\end{thebibliography}
\end{document}
